% 原书第 9 章，印刷页 375--421。
\chapter{多模态大语言模型}
\label{ch:multimodal}
\chaptercredit{9}
\chapterabstract{多模态大语言模型整合文本、图像和音频等多类数据，以模拟人类的感知方式，显著增强人工智能的理解和交互能力。本章介绍多模态大语言模型框架及其组件，结合理论与实际模型，说明每个组件在代表性模型中的不同实现。随后介绍如何把指令调优、上下文学习、思维链和对齐调优等传统大语言模型技术用于多模态情境，以提高跨模态适应性与推理能力。本章详细介绍 Flamingo、Video-LLaMA 和 NExT-GPT 三个当时先进的多模态大语言模型，并说明其各组件如何对应通用框架。在讨论理论基础后，本章以教程展示“文本与其他模态到文本”模型的行为，其中其他模态为图像。教程包括图像分类与描述生成实验，比较零样本、少样本和微调方案，以评估并改善模型表现。}

\section{引言}
人类在现实世界中很少只依赖一种沟通方式。我们通过视觉、听觉和其他感觉输入感知环境，并综合这些信息来理解周围事物并作出反应。多模态大语言模型（multimodal large language model，MMLLM）旨在模拟这种多方面的处理方式，提高模型在实际应用中理解和作出准确响应的能力。此类模型整合文本、图像、音频和其他传感输入等多种数据类型或模态，是人工智能技术的一项重要发展。传统模型只处理单一类型的数据，多模态模型可以处理并解释来自多种模态的复杂数据，输出也可属于不同模态。这模拟了人通过多种感官理解世界并与之互动的认知能力。

跨模态学习涵盖一系列输入与输出涉及不同感知模态的任务，例如视觉和文本数据之间的任务。以下是其中几项。
\begin{enumerate}
\item \textbf{图像与文本检索。}用文本检索相关图像，或用图像检索相关文本描述。
\item \textbf{视频与文本检索。}用文本查找相关视频，或用视频生成文本描述。
\item \textbf{图像和视频描述生成。}为给定图像或视频生成描述性文本。输入是图像或视频等视觉内容，输出是对应的文本描述。
\item \textbf{视觉问答（VQA）。}同时向系统提供图像或视频作为视觉输入，以及以文本形式提出的相关问题。模型需结合视觉和文本输入，根据视觉内容作答。
\item \textbf{带音频反馈的手势控制。}解释手势等视觉输入，并提供相应音频反馈。输入是视觉手势，输出是音频响应或系统所采取的动作，结合了视觉和听觉模态。
\end{enumerate}

\section{简史}
如 Wu et al. (2023c) 所述，随着多模态研究的发展，多模态自动处理领域经历了四个不同的演变阶段。

第一阶段为 1980 至 2000 年，以单模态和统计方法为主。20 世纪 80 年代，统计与图像处理算法广泛用于人脸识别系统的开发。IBM 的研究团队采用隐马尔可夫模型，大幅提高了语音识别的准确性和可靠性（\refcite{ch09-ref-002}{Bahl et al., 1986}）。20 世纪 90 年代，金出武雄的团队率先探索特征脸方法，以主成分分析等统计方法分析人脸图像，有效识别身份（\refcite{ch09-ref-046}{Satoh and Kanade, 1997}）。Dragon Systems 等公司推进了语音识别技术，提高了把口语转换成书面文本的准确率（\refcite{ch09-ref-027}{LaRocca et al., 1999}）。

2000 至 2010 年是第二阶段，其特点是模态转换，尤其重视人机交互。2001 年，AMI 项目探索用计算机记录和处理会议数据，以改善信息检索和协作（\refcite{ch09-ref-005}{Carletta et al., 2005}）。2003 年，认知助理学习与组织项目（Cognitive Assistant that Learns and Organizes，CALO）引入了早期聊天机器人技术，希望创建能够理解并响应人类语言的虚拟助手。它是 Siri 等系统的先驱（\refcite{ch09-ref-053}{Tur et al., 2010}）。社会信号处理（Social Signal Processing，SSP）项目深入分析面部表情、语调等非语言线索，使人机交互更自然（\refcite{ch09-ref-054}{Vinciarelli et al., 2008}）。

2010 至 2020 年是模态融合的第三阶段。深度学习和神经网络融入多模态研究，带来多项突破。2011 年，Ngiam et al. (2011) 提出一种多模态深度学习算法，可联合分析图像和文本等不同模态，改善图像分类、视频分析和语音识别等任务。2012 年，研究人员利用深度玻尔兹曼机捕捉不同模态间的关系，并利用其生成能力（\refcite{ch09-ref-021}{Hinton and Salakhutdinov, 2012}）。2016 年，出现了采用语义注意力的神经图像描述生成算法，可为图像生成描述性文本，从而改善可访问性并支持自动图像标注等应用（\refcite{ch09-ref-065}{You et al., 2016}）。

第四阶段始于 2020 年，其特点是大规模多模态模型的发展，并在本书写作时持续推进。2021 年，对比语言—图像预训练（CLIP）模型改变了以往的做法，以图文配对本身作为学习信号，不依赖预先规定的类别标签（\refcite{ch09-ref-042}{Radford et al., 2021}）。次年，OpenAI 的 DALL-E 2 利用基于 CLIP 图像嵌入的扩散模型，从文本提示生成高质量图像。2023 年，微软发布 KOSMOS-1，这是一种能处理多种模态的信息并通过上下文学习加以适配的多模态大语言模型（\refcite{ch09-ref-023}{Huang et al., 2024}）。PaLM-E 将语言模型与视觉模型相结合，无需针对具体任务微调，在目标检测到代码生成等多种视觉和语言任务上表现出色，成为视觉语言模型的代表（\refcite{ch09-ref-015}{Driess et al., 2023}）。ImageBind 提出为图像、文本、音频、深度、热成像和惯性测量单元（IMU）数据六种模态学习统一的嵌入。它表明只需将其他模态分别与图像配对，就能将这些模态联系起来，用于跨模态检索和生成（\refcite{ch09-ref-018}{Girdhar et al., 2023}）。NExT-GPT 是一个端到端多模态大语言模型，可处理图像、视频、音频和文本输入输出的任意组合（\refcite{ch09-ref-059}{Wu et al., 2023c}）。

\section{多模态大语言模型框架}
\label{sec:ch09-framework}
多模态大语言模型的架构多种多样，取决于其采用的组件，以及为不同功能和模态所作的选择。本节逐一介绍多模态大语言模型的架构组件，并说明每个组件的具体实现方法，见图~\ref{fig:ch09-framework}。该框架综合了多项多模态大语言模型研究，包括 Chip (2023)、Wang et al. (2023)、Wu et al. (2023b)、Xu et al. (2023)、Yin et al. (2023) 和 Zhang et al. (2024a)。

\sourcefigure{fig09_01.pdf}{多模态大语言模型的通用框架，各组件均可采用不同实现。图中图像、音频、视频与三维点云经过模态编码和模态对齐，进入由预训练大语言模型、多模态提示调优、多模态 RLHF 组成的模块。其结果经输出投影器和模态生成器转为所需模态。文本直接送入或由大语言模型输出。}{fig:ch09-framework}

\subsection{模态编码器}
\label{sec:ch09-modality-encoder}
\begin{keypoint}
模态编码器（modality encoder，ME）通常是映射各类数据模态的第一个处理单元。一般来说，图像、视频或音频等每类数据都经过相应的专用编码器。这些编码器将不同数据类型的特有属性转化为嵌入，即模型后续各层可以统一理解和处理的向量表示。
\end{keypoint}

模态编码器的运算可表示为
\begin{equation}
F_X=\mathrm{ME}_X(I_X),
\label{eq:ch09-modality-encoder}
\end{equation}
其中，$I_X$ 表示图像、视频、音频或三维物体等模态的输入数据，$F_X$ 表示提取的特征。下面讨论研究中用于不同数据流的若干常见编码器。

Dosovitskiy et al. (2020) 提出的视觉 Transformer（ViT）把此前常用于自然语言处理的 Transformer 架构用于图像分析。ViT 将图像划分为图块并作线性投影，再用 Transformer 块对视觉信息编码。基于 ViT，Radford et al. (2021) 的 CLIP ViT 利用自然语言监督学习视觉概念。CLIP ViT 在大量图文配对数据上进行对比学习，改善图像与相应文本描述之间的对齐，大幅提高模型理解视觉内容并分类的能力。许多多模态大语言模型使用 CLIP 编码器处理图像数据。Fang et al. (2023) 提出的 EVA-CLIP ViT 进一步改进了这一方法，解决了训练 CLIP 等大规模模型时遇到的部分问题。

\begin{practicaltip}
EVA-CLIP ViT 改善了训练过程的稳定性，可高效扩展模型规模，并改善视觉识别任务中多模态基础模型的训练，因而适合作为图像编码器。
\end{practicaltip}

处理视频时，常见做法是采样固定数量的帧（通常为五帧），再采用与图像相同的预处理步骤。统一处理使不同视觉模态的特征提取方式保持一致。

音频领域已有多种用于变换声音数据的编码器，如 C-Former、HuBERT、BEATs 和 Whisper。Chen et al. (2023) 的 C-Former 结合连续积分发放（continuous integrate-and-fire，CIF）对齐机制和 Transformer 架构来执行序列转换，从原始声音数据中提取细致的音频特征。Shi et al. (2022) 提出的 HuBERT 采用源自 BERT 框架的自监督学习方法，预测被掩码的音频片段，以学习可服务于不同下游任务的语音表示。Chen et al. (2022) 的 BEATs 则采用迭代式音频预训练框架。

三维视觉理解可用于增强现实与虚拟现实（AR 与 VR）、自动驾驶导航、元宇宙和机器人等领域，因而受到重视。继 ULIP 框架之后，ULIP-2 为三维物体自动生成描述语言，在不依赖人工标注的情况下构建大规模三模态数据集，成为三维视觉理解的代表性方法。三维点云模态可用 ULIP-2 模型的编码器编码（\refcite{ch09-ref-062}{Xue et al., 2023}）。

\subsection{输入投影器}
\label{sec:ch09-input-projector}
\begin{keypoint}
模态对齐 $\mathrm{IN\_ALIGN}_{X\to T}$ 的核心是使用输入投影器，将不同模态的编码特征 $F_X$ 对齐到文本特征空间 $T$。对齐后得到提示 $P_X$，与文本特征 $F_T$ 一起输入大语言模型。
\end{keypoint}

\transnote{此处原文把待投影的特征也写作 $F_T$。据式~\eqref{eq:ch09-input-projection}，非文本输入特征应记为 $F_X$。}

给定多模态与文本的配对数据集 $\{(I_X,t)\}$，主要目标是最小化在模态 $X$ 条件下生成文本的损失 $L_{\text{txt-gen}}$。
\begin{equation}
\arg\min_{\mathrm{IN\_ALIGN}_{X\to T}}
L_{\text{txt-gen}}\bigl(\mathrm{LLM}(P_X,F_T),t\bigr).
\label{eq:ch09-input-loss}
\end{equation}
其中，对非文本特征应用输入投影器，就得到作为提示的对齐特征 $P_X$。
\begin{equation}
P_X=\mathrm{IN\_ALIGN}_{X\to T}(F_X).
\label{eq:ch09-input-projection}
\end{equation}

多模态预训练模型用多层 Transformer 架构提取各模态特征，并让特征相互作用。按照整合多模态信息的方式，可将这些架构分为单流和跨流两类。
\begin{itemize}
\item \textbf{单流架构。}将图像、文本等多模态输入同等对待，在统一模型中融合。先从各模态提取单模态特征，再将特征表示为词元并用分隔符拼接，如图~\ref{fig:ch09-single-stream}所示。拼接后的特征输入多模态 Transformer。多头自注意力机制让各单模态特征交互并融合，生成多模态融合特征（\refcite{ch09-ref-034}{Li et al., 2020c}）。通常用 Transformer 的分类词元取得这类特征，该词元汇集了各模态的信息，提高模型的表征能力。
\item \textbf{跨流架构。}不同模态的特征由独立模型并行提取，再通过下文介绍的自监督对比学习完成对齐，如图~\ref{fig:ch09-cross-stream}所示。跨流架构主要对齐各个单模态特征，单流架构则产生多模态融合特征。\footnote{原文该句的关系从句误把单流架构说成着重对齐单模态特征，此处按相邻两段及图~\ref{fig:ch09-single-stream}、图~\ref{fig:ch09-cross-stream}所述的过程译写。}采用跨流结构的大规模多模态语言模型有 BriVL 与 CLIP（\refcite{ch09-ref-025}{Huo et al., 2021}; \refcite{ch09-ref-042}{Radford et al., 2021}）。它们的一个重要特点是将不同模态的特征对齐到同一个高维特征空间。跨流模型很灵活，改变一种模态的结构不会影响其他模态，因而更容易用于实际应用。这些模型主要用于嵌入层面的匹配，通常利用所谓的“弱语义相关性”，并考虑多模态数据之间的差异与互补关系。此类预训练模型的结构还取决于具体的预训练目标。执行重建掩码图像或根据文本描述生成图像等生成任务时，模型会在编码器后接解码器，把多模态融合特征转换成相应输出，以完成预训练。
\end{itemize}

\sourcefigure{fig09_02.pdf}{单流架构。将文本、图像和音频分别编码、嵌入并词元化，拼成同一序列，经多模态 Transformer 融合以执行预训练任务。}{fig:ch09-single-stream}
\sourcefigure{fig09_03.pdf}{跨流架构。不同模态由独立编码器产生嵌入，并用对比学习对齐。图中也给出了将嵌入拼接送入多模态 Transformer 的路径。}{fig:ch09-cross-stream}

多模态 Transformer 通过自注意力机制及其变体实现融合、对齐等跨模态交互。自注意力可以处理来自任意模态的词元嵌入，不受模态、词元化方法和嵌入方式的限制。给定来自两种不同模态的输入 $X_A$ 和 $X_B$，用 $Z^{(A)}$ 和 $Z^{(B)}$ 表示各自的词元嵌入。下面以两种模态为例说明各种方法及其数学形式。这些方法可推广到更多模态。

\begin{enumerate}
\item \textbf{早期求和。}在 Transformer 各层处理之前，将不同模态的词元嵌入在每个位置加权求和。
\[
Z\leftarrow T_f\bigl(\alpha Z^{(A)}\oplus\beta Z^{(B)}\bigr)
=\mathrm{MHSA}\bigl(Q^{(AB)},K^{(AB)},V^{(AB)}\bigr),
\]
其中，$\oplus$ 表示逐元素求和。这种方法简单有效，无需增加计算复杂度（\refcite{ch09-ref-017}{Gavrilyuk et al., 2020}）。
\item \textbf{早期拼接（Co-Transformer）。}将不同模态的词元嵌入序列拼接起来。
\[
Z\leftarrow T_f\bigl(C(Z^{(A)},Z^{(B)})\bigr).
\]
这种全注意力或 Co-Transformer 方法将拼接结果作为一个序列处理，用其他模态提供的上下文增强各模态的编码（\refcite{ch09-ref-048}{Sun et al., 2019}）。
\item \textbf{分层注意力（多流到单流）。}先由各自的 Transformer 流编码多模态输入，再拼接并融合它们的输出。
\[
Z\leftarrow T_{f3}\bigl(C(T_{f1}(Z^{(A)}),T_{f2}(Z^{(B)}))\bigr).
\]
该方法是一种后期交互或融合，也可视为早期拼接的特例（\refcite{ch09-ref-032}{Li et al., 2021}）。
\item \textbf{分层注意力（单流到多流）。}先用共享的单流 Transformer 编码拼接后的多模态输入，再用各自的流对每个模态编码。
\[
\begin{aligned}
C(Z^{(A)},Z^{(B)})&\leftarrow T_{f1}\bigl(C(Z^{(A)},Z^{(B)})\bigr),\\
Z^{(A)}&\leftarrow T_{f2}(Z^{(A)}),\\
Z^{(B)}&\leftarrow T_{f3}(Z^{(B)}).
\end{aligned}
\]
InterBERT 采用此结构捕捉跨模态交互，同时保留单模态表征的独立性（\refcite{ch09-ref-036}{Lin et al., 2020}）。
\item \textbf{交叉注意力。}双流 Transformer 交换各流的查询嵌入，以增强跨模态交互。
\[
\begin{aligned}
Z^{(A)}&\leftarrow\mathrm{MHSA}(Q_B,K_A,V_A),\\
Z^{(B)}&\leftarrow\mathrm{MHSA}(Q_A,K_B,V_B).
\end{aligned}
\]
这种方法首先在 ViLBERT 中提出，兼顾计算效率和跨模态感知（\refcite{ch09-ref-039}{Lu et al., 2019}）。
\item \textbf{交叉注意力后拼接。}将交叉注意力的各流拼接后进一步处理，以建模全局上下文。
\[
\begin{aligned}
Z^{(A)}&\leftarrow\mathrm{MHSA}(Q_B,K_A,V_A),\\
Z^{(B)}&\leftarrow\mathrm{MHSA}(Q_A,K_B,V_B),\\
Z&\leftarrow T_f\bigl(C(Z^{(A)},Z^{(B)})\bigr).
\end{aligned}
\]
这种分层跨模态交互减轻了单用交叉注意力的缺点（\refcite{ch09-ref-069}{Zhan et al., 2021}）。
\end{enumerate}

% 原书第 9.3.3 节。
\subsection{预训练：核心大语言模型、数据集和任务目标}
\label{sec:ch09-pretraining}
生成响应的大语言模型是多模态大语言模型的核心。输入既可能包含文本数据，也可能包含非文本数据，因而需要新的技术，让语言模型以多种模态为条件生成响应。大语言模型处理来自各模态的表征，以理解输入语义、进行推理和作出决策。它产生两种主要输出。
\begin{enumerate}
\item 直接输出的文本，记为 $t$。
\item 其他模态的信号词元，记为 $S_X$。
\end{enumerate}
这些信号词元作为指令，告诉生成器是否要生成多模态内容。其数学表示为
\begin{equation}
(t,S_X)=\mathrm{LLM}(P_X,F_T),
\label{eq:ch09-mmlm-output}
\end{equation}
其中，$P_X$ 可视为对大语言模型进行软提示调优所用的提示。

多模态大语言模型可以分为纯编码器、纯解码器和编码器—解码器三类。表~\ref{tab:ch09-base-models}列出常用于多模态训练的一些大语言模型及相关细节。

{\small\setlength{\tabcolsep}{4pt}\renewcommand{\arraystretch}{1.13}
\begin{longtable}{@{}>{\raggedright\arraybackslash}p{31mm} >{\raggedright\arraybackslash}p{22mm} >{\raggedright\arraybackslash}p{56mm}@{}}
\caption{多模态系统采用的基础大语言模型。}\label{tab:ch09-base-models}\\
\toprule
大语言模型 & 架构 & 说明\\
\midrule
\endfirsthead
\multicolumn{3}{@{}l}{表 \thetable\ （续）}\\
\toprule
大语言模型 & 架构 & 说明\\
\midrule
\endhead
\midrule
\multicolumn{3}{r@{}}{下页续}\\
\endfoot
\bottomrule
\endlastfoot
Flan-T5（\refcite{ch09-ref-012}{Chung et al., 2022}） & 编码器—解码器 & 探索对 T5 进行指令调优，展示了较强的零样本学习和思维链（CoT）能力。\\
ChatGLM2（\refcite{ch09-ref-068}{Zeng et al., 2022}） & 自回归 & 中英双语对话模型，针对中文问答与对话优化。\\
UL2（\refcite{ch09-ref-050}{Tay et al., 2022}） & 编码器—解码器 & 用去噪目标训练，在基准上超过 T5。\\
Qwen（\refcite{ch09-ref-003}{Bai et al., 2023}） & 纯解码器 & 侧重中英双语能力，采用对齐技术改善对话表现。\\
Chinchilla（\refcite{ch09-ref-022}{Hoffmann et al., 2022}） & 因果解码器 & 主张模型规模随数据集规模扩展，在大规模文本语料上训练。\\
OPT（\refcite{ch09-ref-074}{Zhang et al., 2022}） & 因果解码器 & 尝试以开源方式复现 GPT-3 的性能。\\
PaLM（\refcite{ch09-ref-011}{Chowdhery et al., 2022}） & 因果解码器 & 注意力层和前馈层并行，采用 RoPE 嵌入、SwiGLU 激活函数等方式加快训练。\\
Llama（\refcite{ch09-ref-052}{Touvron et al., 2023}） & 纯解码器 & 在纯解码器架构中使用高效的因果注意力。\\
Llama-2（\refcite{ch09-ref-052}{Touvron et al., 2023}） & 纯解码器 & 相对 Llama 增加 40\% 训练数据，并改进了对话生成。\\
Vicuna（\refcite{ch09-ref-009}{Chiang et al., 2023}） & 纯解码器 & 以 Llama 为基础，利用用户对话数据训练，以增强对话能力。\\
\end{longtable}
}

在预训练阶段，模型通常使用包含图文、视频与文本、音频与文本等多种模态的数据集。这一阶段的重点是训练输入投影器与输出投影器，使不同模态的特征对齐。优化一般集中于这两个组件，有时也会对大语言模型进行参数高效微调，进一步改善模型处理多模态信息的能力。

表~\ref{tab:ch09-pretraining-data}列出预训练中常用的数据集（\refcite{ch09-ref-055}{Wang et al., 2023}; \refcite{ch09-ref-063}{Yin et al., 2023}）。表中 K、M 和 B 分别表示千、百万和十亿。

{\small\renewcommand{\arraystretch}{1.04}
\begin{longtable}{@{}p{49mm}p{30mm}l@{}}
\caption{预训练中常用的数据集。}\label{tab:ch09-pretraining-data}\\
\toprule
数据集 & 模态 & 规模\\
\midrule
\endfirsthead
\multicolumn{3}{@{}l}{表 \thetable\ （续）}\\
\toprule
数据集 & 模态 & 规模\\
\midrule
\endhead
\midrule
\multicolumn{3}{r@{}}{下页续}\\
\endfoot
\bottomrule
\endlastfoot
ALIGN & 图像—文本 & 1.8B\\
LTIP & 图像—文本 & 312M\\
MS-COCO & 图像—文本 & 620K\\
VisualGenome & 图像—文本 & 4.5M\\
CC3M & 图像—文本 & 3.3M\\
CC12M & 图像—文本 & 12.4M\\
SBU & 图像—文本 & 1M\\
LAION-400M & 图像—文本 & 400M\\
Flickr30k & 图像—文本 & 158K\\
AIChallengerCaptions & 图像—文本 & 1.5M\\
COYO & 图像—文本 & 747M\\
Wukong & 图像—文本 & 101M\\
COCOCaption & 图像—文本 & 1M\\
WebLI & 图像—文本 & 12B\\
EpisodicWebLI & 图像—文本 & 400M\\
CC595k & 图像—文本 & 595K\\
RefCOCO+ & 图像—文本 & 142K\\
Visual-7W & 图像—文本 & 328K\\
OCR-VQA & 图像—文本 & 1M\\
ST-VQA & 图像—文本 & 32K\\
DocVQA & 图像—文本 & 50K\\
TextVQA & 图像—文本 & 45.3K\\
DataComp & 图像—文本 & 1.4B\\
GQA & 图像—文本 & 22M\\
VQAv2 & 图像—文本 & 1.4M\\
DVQA & 图像—文本 & 3.5M\\
OK-VQA & 图像—文本 & 14K\\
A-OKVQA & 图像—文本 & 24.9K\\
TextCaptions & 图像—文本 & 145K\\
M3W（交错） & 图像—文本 & 43.3M 个实例\\
MMC4（交错） & 图像—文本 & 101.2M 个实例\\
MSRVTT & 视频—文本 & 200K\\
WebVid & 视频—文本 & 10M\\
VTP & 视频—文本 & 27M\\
AISHELL-2 & 音频 & 128K\\
WaveCaps & 音频 & 403K\\
VSDial-CN & 图像—音频—文本 & 1.2M\\
\end{longtable}
}

针对任务和模态设计学习目标，是多模态预训练的重要工作。下面介绍预训练中常用的学习目标。

\subsubsection{对比学习}
CLIP 出现之前，视觉语言模型主要采用分类或语言建模目标。分类方法受限于预定义类别，模型的响应形式和对不同任务的适应性也受限制。语言建模目标更灵活，但着重于为每张图像生成特定文本，训练有一定难度。

CLIP 的对比学习试图克服以往模型的局限，将目标从预测每张图像对应的确切文本，转为判断给定文本与某张图像的关联是否强于其他图像（\refcite{ch09-ref-042}{Radford et al., 2021}）。对于一个包含 $N$ 对图文样本的批次，CLIP 生成 $N$ 个文本嵌入和 $N$ 个图像嵌入。记 $V_1,V_2,\ldots,V_N$ 为这 $N$ 张图像的嵌入，$L_1,L_2,\ldots,L_N$ 为这 $N$ 段文本的嵌入。CLIP 计算 $V_i$ 和 $L_j$ 的 $N^2$ 种可能配对的余弦相似度。训练目标是最大化 $N$ 个正确配对的相似度，同时最小化 $N^2-N$ 个错误配对的相似度。
\begin{align}
L_{\mathrm{i2t}}&=-\frac{1}{N}\sum_i
\log\frac{\exp(V_i^{T}L_i/\sigma)}{\sum_j\exp(V_i^{T}L_j/\sigma)},
\label{eq:ch09-image-to-text-loss}\\
L_{\mathrm{t2i}}&=-\frac{1}{N}\sum_i
\log\frac{\exp(L_i^{T}V_i/\sigma)}{\sum_j\exp(L_i^{T}V_j/\sigma)},
\label{eq:ch09-text-to-image-loss}\\
L_{\mathrm{CL}}&=L_{\mathrm{i2t}}+L_{\mathrm{t2i}}.
\label{eq:ch09-contrastive-loss}
\end{align}
这里 $L_{\mathrm{i2t}}$ 和 $L_{\mathrm{t2i}}$ 分别为图像到文本与文本到图像的分类损失，$L_{\mathrm{CL}}$ 为总对比损失。$V_i$ 和 $L_i$ 分别表示归一化的图像嵌入与文本嵌入，$N$ 为批大小，$\sigma$ 为温度参数。

\subsubsection{模态匹配损失}
模态匹配损失（modality matching loss，MML）在大型多模态模型的预训练中起重要作用，可以捕捉不同模态间显式或隐式的对齐关系。Unicoder-VL 等模型采用这种损失函数，通过视觉语言匹配（visual linguistic matching，VLM）进行视觉语言预训练（\refcite{ch09-ref-029}{Li et al., 2020a}）。VLM 方法抽取图像与句子的正、负配对，训练模型识别配对是否对齐。其目标是预测给定样本对的匹配分数。
\begin{equation}
\begin{aligned}
L_{\mathrm{MML}}={}&-\sum_{(x,y)\in\mathrm{Pos}}
  \log p(\text{对齐}\mid x,y)\\
&-\sum_{(x',y')\in\mathrm{Neg}}
  \log p(\text{未对齐}\mid x',y').
\end{aligned}
\label{eq:ch09-modality-matching-loss}
\end{equation}
其中 $(x,y)$ 表示图像与句子的正配对，$(x',y')$ 表示负配对。模型预测配对已经对齐的概率 $p(\text{对齐}\mid x,y)$ 和未对齐的概率 $p(\text{未对齐}\mid x',y')$。

InterBERT 提出了采用困难负例的图文匹配变体 ITM-hn（\refcite{ch09-ref-036}{Lin et al., 2020}）。它从 TF-IDF 相似度最高的样本中选择负例，有别于常规的负例采样策略。
\begin{equation}
\begin{aligned}
L_{\mathrm{ITM\text{-}hn}}={}&-\sum_{(x,y)\in\mathrm{Pos}}
  \log p(\text{对齐}\mid x,y)\\
&-\sum_{(x',y')\in\text{Hard Neg}}
  \log p(\text{未对齐}\mid x',y').
\end{aligned}
\label{eq:ch09-hard-negative-matching-loss}
\end{equation}
纳入 TF-IDF 相似度较高的困难负例后，模型必须区分关系紧密但没有正确对齐的配对，学习因而更困难，也更有效。

\subsubsection{掩码语言建模}
掩码语言建模（MLM）是预训练框架中的常用目标。研究人员通常随机用特殊词元掩盖输入词，并要求模型补全被掩盖的词。模型利用周围的词和相关图像区域作为上下文。Wang et al. (2021) 开发的 SIMVLM 将这种方法与前缀语言建模（PrefixLM）结合。PrefixLM 对前缀序列施加双向注意力，并对后续词元采用自回归分解。这里将词记作 $w=\{x_1,\ldots,x_K\}$，图像区域记作 $v=\{v_1,\ldots,v_T\}$。在 MLM 中，随机生成一组索引 $m$，掩盖占比为 $p\%$ 的输入词 $x_m$。其目标是利用未掩盖的词 $x_{\neg m}$ 和全部图像区域 $v$ 预测这些词，通过最小化负对数似然来实现。
\begin{equation}
L_{\mathrm{MLM}}(\theta)=-\mathbb{E}_{(x,v)}
\log P_\theta(x_m\mid x_{\neg m},v),
\label{eq:ch09-mlm-loss}
\end{equation}
其中 $\theta$ 为可训练参数。

除了 MLM，SIMVLM 中的 PrefixLM 也是一种预训练视觉语言表征的策略。给定前缀，其目标是预测文本序列的后续部分。
\begin{equation}
L_{\mathrm{PrefixLM}}(\theta)=-\mathbb{E}_{x\sim D}
\log P_\theta(x_{\geq T_p}\mid x_{<T_p}),
\label{eq:ch09-prefixlm-loss}
\end{equation}
其中 $x$ 是文本序列，$D$ 表示预训练数据，$T_p$ 是前缀词元序列的长度。

\subsubsection{掩码物体分类}
这种技术有选择地掩盖图像的部分区域，通常将相应特征值置零，然后将物体检测器预测的标签用作被掩盖区域的真实标签。

掩码物体分类（MOC）与自然语言处理中的 MLM 方法相似。MOC 以一定概率 $p\%$ 修改图像区域的视觉特征，以遮蔽特定区域。主要目标是正确预测被掩盖的图像区域 $v_m^i$ 中的物体类别。将这些区域对应的编码器输出送入全连接层（FC），计算 $T$ 个物体类别的分数（\refcite{ch09-ref-029}{Li et al., 2020a}）。再用 softmax 函数将分数转化为归一化分布 $g_\theta(v_m^i)$。MOC 目标形式化为
\begin{equation}
L_{\mathrm{MOC}}(\theta)=-\mathbb{E}_{(w,v)}
\left[\sum_{i=1}^M \mathrm{CE}\bigl(c(v_m^i),g_\theta(v_m^i)\bigr)\right],
\label{eq:ch09-moc-loss}
\end{equation}
其中 $c(v_m^i)$ 是被掩盖图像区域的真实标签，$\mathrm{CE}$ 为交叉熵损失函数，$\theta$ 为模型参数，期望 $\mathbb{E}$ 对词 $w$ 和视觉特征 $v$ 的分布求取。MOC 旨在增强模型在视觉内容仅部分可见或被遮挡时推断并分类物体的能力，从而提高其理解视觉信息的能力。\footnote{式~\eqref{eq:ch09-moc-loss}保留底本的负号。若 $\mathrm{CE}$ 按通常定义表示非负交叉熵，则最小化交叉熵时不再额外取负号。}

\subsubsection{图文匹配}
图文匹配（image-text matching，ITM）是开发能理解视觉内容及其文本描述之间关系的模型的重要环节。ITM 需要生成负训练数据，通常为每张图像配上不相关的句子，或为每个句子配上不相关的图像。目标是提高模型区分正确配对与错误配对的能力。

每个图文对 $(v,t)$ 带有一个真实标签 $y$，表示配对是否正确，即是正例还是负例。ITM 使用二分类损失函数优化模型，衡量模型能否正确判断这种对齐。其损失函数 $L_{\mathrm{ITM}}(\theta)$ 为
\begin{equation}
\begin{aligned}
L_{\mathrm{ITM}}(\theta)=-\mathbb{E}_{(v,t)}
\bigl[&y\log s_\theta(v,t)\\
&+(1-y)\log\bigl(1-s_\theta(v,t)\bigr)\bigr],
\end{aligned}
\label{eq:ch09-itm-loss}
\end{equation}
其中 $s_\theta(v,t)$ 是参数为 $\theta$ 的模型算得的图文相似度分数。期望 $\mathbb{E}_{(v,t)}$ 对图文对的分布求取。此损失函数衡量模型区分正确与错误配对的能力，以改善对视觉与文本模态之间复杂关系的理解。

\subsubsection{图像到文本的生成}
图像到文本的生成（image-text generation，ITG）是视觉语言预训练的一项重要任务，利用已对齐的图文对训练模型，让它能根据给定图像生成文本。例如，Xu et al. (2021) 用 ITG 目标训练 E2E-VLP 模型。ITG 目标表示为
\begin{equation}
L_{\mathrm{ITG}}=-\sum_{(x,y)\in(X,Y)}
\log\prod_{t=1}^{n}P(y_t\mid y_{<t},x).
\label{eq:ch09-itg-loss}
\end{equation}
其中 $X$ 表示带有上下文的视觉序列，$Y$ 表示生成文本的集合，$n$ 为文本 $y$ 的词元数。该目标旨在最大化正确生成文本词元序列的概率。每个词元 $y_t$ 的生成以前序词元 $y_{<t}$ 和视觉输入 $x$ 为条件。

\subsubsection{视频与字幕匹配}
视频与字幕匹配（video-subtitle matching，VSM）可用于视频与文本的预训练。以 HERO 为例，它考虑局部与全局对齐两个目标（\refcite{ch09-ref-031}{Li et al., 2020b}）。评分函数分别衡量视频与字幕在局部和全局上的对齐程度。优化时，模型要分辨正确匹配的视频与字幕对（正例）和错误匹配的配对（负例）。

在 HERO 中，局部对齐分数函数与全局对齐分数函数为
\[
\begin{aligned}
S_{\mathrm{local}}(s_q,v)&=V^{\mathrm{temp}}q\in\mathbb{R}^{N_v},\\
S_{\mathrm{global}}(s_q,v)&=\max\left(
  \frac{V^{\mathrm{temp}}}{\lVert V^{\mathrm{temp}}\rVert}
  \cdot\frac{q}{\lVert q\rVert}
\right).
\end{aligned}
\]
正负查询—视频对的铰链损失、局部对齐损失与全局对齐损失为
\[
\begin{aligned}
L_h(S_{\mathrm{pos}},S_{\mathrm{neg}})
 &=\max(0,\delta+S_{\mathrm{neg}}-S_{\mathrm{pos}}),\\
L_{\mathrm{local}}
 &=-\mathbb{E}_D\bigl[\log p_{\mathrm{st}}[y_{\mathrm{st}}]
                      +\log p_{\mathrm{ed}}[y_{\mathrm{ed}}]\bigr],\\
L_{\mathrm{global}}
 &=\mathbb{E}_D\bigl[
    L_h\bigl(S_{\mathrm{global}}(s_q,v),S_{\mathrm{global}}(\hat{s}_q,v)\bigr)\\
 &\hspace{4.9em}{}+L_h\bigl(S_{\mathrm{global}}(s_q,v),S_{\mathrm{global}}(s_q,\hat{v})\bigr)
    \bigr].
\end{aligned}
\]
合并得到 VSM 损失
\[
L_{\mathrm{VSM}}=\lambda_1 L_{\mathrm{local}}+\lambda_2 L_{\mathrm{global}}.
\]
\transnote{底本把铰链项写成 $\delta+S_{\mathrm{pos}}-S_{\mathrm{neg}}$，在全局损失的期望前加负号，局部损失的第一个对数括号也未闭合。此处依 \href{https://github.com/linjieli222/HERO/blob/master/model/pretrain.py}{HERO 作者的实现}排式，即局部损失是起点与终点交叉熵之和，全局损失最小化正例分数未高出负例分数至少 $\delta$ 时的惩罚。}

在该模型中，$s_q$ 是从所有字幕句子中抽取的查询，$v$ 是完整视频片段，$V^{\mathrm{temp}}\in\mathbb{R}^{N_v\times d}$ 是时间 Transformer 最终生成的视频帧表征。查询向量 $q\in\mathbb{R}^d$，起点和终点索引 $y_{\mathrm{st}},y_{\mathrm{ed}}\in\{1,\ldots,N_v\}$，概率向量 $p_{\mathrm{st}},p_{\mathrm{ed}}\in\mathbb{R}^{N_v}$ 来自对齐分数。铰链损失 $L_h$ 比较正例与负例的查询—视频对，其中 $(s_q,v)$ 为正例，$(s_q,\hat{v})$ 与 $(\hat{s}_q,v)$ 为负例。间隔超参数 $\delta$ 和权重系数 $\lambda_1,\lambda_2$ 用于控制这套方法。

\subsubsection{帧顺序建模}
在 HERO 中，帧顺序建模（frame order modeling，FOM）被视为一个分类问题，目标是正确预测一组给定视频帧的时间顺序（\refcite{ch09-ref-031}{Li et al., 2020b}）。FOM 主要用于确定从视频中抽出的一部分帧的原始时间戳序列，以考察模型理解视频时间变化和叙事顺序的能力。

FOM 目标表示为损失函数
\begin{equation}
L_{\mathrm{FOM}}=-\mathbb{E}\left[\sum_{i=1}^{R}\log P[r_i,t_i]\right],
\label{eq:ch09-frame-order-loss}
\end{equation}
其中各符号含义如下。
\begin{itemize}
\item $R$ 表示顺序被打乱、需要分类的帧的总数。
\item $i$ 是这组帧中的索引，取值为 $1$ 至 $R$。
\item $t_i$ 表示第 $i$ 帧在视频中的真实时间戳位置，取值为 $1$ 至 $N_v$，其中 $N_v$ 为视频的总帧数。
\item $r_i$ 是第 $i$ 帧在重新排序后的位置索引。
\item $P$ 是 $N_v\times N_v$ 的概率矩阵，元素 $P[r_i,t_i]$ 表示模型预测重排后位置 $r_i$ 的帧对应时间戳 $t_i$ 的概率。
\end{itemize}

% 原书第 9.3.4 至 9.3.7 节。
\subsection{多模态大语言模型的调优与增强}
预训练完成后，可以进一步提高多模态大语言模型的适应性、推理和任务泛化能力。这里介绍三种方法。多模态指令调优（MM-IT）训练模型遵循多种任务指令。多模态上下文学习（MM-ICL）使模型能用已有知识完成输入提示中的新任务。多模态思维链（MM-CoT）使模型求解复杂问题的过程更透明，也更有条理。

\subsubsection{多模态指令调优}
\label{sec:ch09-mm-it}
\begin{keypoint}
传统的监督微调需要大量数据，少样本提示方法带来的提升也有限。指令调优（IT）旨在把任务能力推广到初始训练数据之外，参见第~4.2~节。多模态指令调优（MM-IT）沿用指令调优的原理，在按指令任务组织的多模态数据集上微调大语言模型（\refcite{ch09-ref-038}{Liu et al., 2024}; \refcite{ch09-ref-076}{Zhao et al., 2023}; \refcite{ch09-ref-077}{Zhu et al., 2023}）。模型能通过理解指令处理新任务，从而提高在不同模态上的零样本学习能力。
\end{keypoint}

多模态指令样本可表示为三元组 $(I,M,R)$，分别对应指令、多模态输入和真实响应。参数为 $\theta$ 的模型根据指令和多模态输入预测答案。
\begin{equation}
A=f(I,M;\theta).
\label{eq:ch09-mm-it-prediction}
\end{equation}
这里 $A$ 表示预测答案。训练常沿用原先的自回归目标，要求多模态大语言模型预测响应中的下一个词元。
\begin{equation}
L(\theta)=-\sum_{i=1}^N\log p(R_i\mid I,R_{<i};\theta),
\label{eq:ch09-mm-it-loss}
\end{equation}
其中 $N$ 为真实响应长度。模型根据前文与指令准确生成响应序列中的下一个词元。\footnote{式~\eqref{eq:ch09-mm-it-loss}保留原书写法。按式~\eqref{eq:ch09-mm-it-prediction}，多模态输入 $M$ 也属于生成响应的条件。}图~\ref{fig:ch09-mm-it}展示视觉问答任务的一个示例模板，表~\ref{tab:ch09-mm-it}列出若干常用的多模态指令调优数据集。

\sourcefigure{fig09_04.pdf}{视觉问答任务的多模态指令调优模板。在序列开始标记 \texttt{<BOS>} 后，要求按给定的任务指令作答。模板包含指令、由图像与文本组成的输入、响应，以及序列结束标记 \texttt{<EOS>}。}{fig:ch09-mm-it}

{\small\setlength{\tabcolsep}{4pt}
\begin{longtable}{@{}p{41mm}p{43mm}r@{}}
\caption{多模态指令调优数据集。第二列表示输入到输出的模态。T 为文本，I 为图像，V 为视频，A 为音频，B 为边界框，3D 为点云。}\label{tab:ch09-mm-it}\\
\toprule
数据集 & 输入到输出 & 实例数\\
\midrule
\endfirsthead
\multicolumn{3}{@{}l}{表 \thetable\ （续）}\\
\toprule
数据集 & 输入到输出 & 实例数\\
\midrule
\endhead
\midrule
\multicolumn{3}{r@{}}{下页续}\\
\endfoot
\bottomrule
\endlastfoot
MiniGPT-4’sIT & I+T$\to$T & 5K\\
StableLLaVA & I+T$\to$T & 126K\\
LLaVA’sIT & I+T$\to$T & 150K\\
SVIT & I+T$\to$T & 3.2M\\
LLaVAR & I+T$\to$T & 174K\\
ShareGPT4V & I+T$\to$T & —\\
DRESS’sIT & I+T$\to$T & —\\
VideoChat’sIT & V+T$\to$T & 11K\\
Video-ChatGPT’sIT & V+T$\to$T & 100K\\
Video-LLaMA’sIT & I/V+T$\to$T & 171K\\
InstructBLIP’sIT & I/V+T$\to$T & $\sim$1.6M\\
X-InstructBLIP’sIT & I/V/A/3D+T$\to$T & $\sim$1.8M\\
MIMIC-IT & I/V+T$\to$T & 2.8M\\
PandaGPT’sIT & I+T$\to$T & 160K\\
MGVLID & I+B+T$\to$T & —\\
M3IT & I/V/B+T$\to$T & 2.4M\\
LAMM & I+3D+T$\to$T & 196K\\
BuboGPT’sIT & (I+A)/A+T$\to$T & 9K\\
T2M & T$\to$I/V/A+T & 14.7K\\
MosIT & I+V+A+T$\to$I+V+A+T & 5K\\
\end{longtable}
}

\subsubsection{多模态上下文学习}
\label{sec:ch09-mm-icl}
\begin{keypoint}
上下文学习（ICL）使大语言模型能从少数示例中理解任务并加以执行，提示通常还包含指令。传统监督学习需要大量数据，上下文学习则无需额外训练就能处理新任务。指令调优需要在指令数据集上微调模型，上下文学习在推理时借助模型已预训练的能力适应新任务，无需更新模型。

把上下文学习扩展到多模态，就得到多模态上下文学习（MM-ICL），它用多种模态丰富学习过程（\refcite{ch09-ref-019}{Gupta and Kembhavi, 2023}）。推理时，MM-ICL 将一组演示样例与原待处理样本同时输入，用多个上下文示例丰富模型可以利用的信息。
\end{keypoint}

\sourcefigure{fig09_05.pdf}{图像描述生成任务中的多模态上下文学习。给出“生成下列图像的描述”这一指令后，示范如何描述坐在沙发上的猫和在山间小径远足的一群人。最后输入一张房前停着汽车的图像，留出响应让模型生成。}{fig:ch09-mm-icl}

图~\ref{fig:ch09-mm-icl}用两个示例展示如何用 MM-ICL 生成图像描述。示例的结构和数量可以灵活调整，模型表现往往也取决于它们的排列顺序。表~\ref{tab:ch09-mm-icl}列出两个重要的 MM-ICL 数据集。

\begin{table}[!htbp]
\centering\small
\caption{多模态上下文学习数据集。}
\label{tab:ch09-mm-icl}
\setlength{\tabcolsep}{3pt}
\begin{tabularx}{\linewidth}{@{}llr>{\raggedright\arraybackslash}X@{}}
\toprule
数据集 & 模态 & 规模 & 说明\\
\midrule
MM-ICL & 图像—文本 & 5.8M & 包含交错的文本与图像输入，以及人工构造的多模态上下文学习输入。\\
MIMIC-IT & 图像—文本 & 2.8M & 提供多模态指令与响应配对，提高视觉语言模型在多种语言中的感知、推理和规划能力。\\
\bottomrule
\end{tabularx}
\end{table}

\subsubsection{多模态思维链}
\label{sec:ch09-mm-cot}
\begin{keypoint}
多模态思维链（multimodal chain-of-thought，MM-CoT）是大语言模型思维链概念的扩展。思维链在复杂推理任务中有效，要求大语言模型同时生成最终答案和中间推理步骤，类似人类的认知过程。MM-CoT 将单模态思维链用于多模态情境，首先需要在不同模态之间建立联系。可以融合特征，也可以把视觉输入转为文本描述。MM-CoT 可以通过微调建立，也可采用无需训练的少样本或零样本学习，各方法所需的样本量不同。
\end{keypoint}

Lian et al. (2023) 使用 ChatGPT 将人工标注者提供的多条描述线索综合为连贯的摘要，提取其中的主要行为和表情，再用这些综合线索推断人物的实际情绪，如图~\ref{fig:ch09-emotion-cot}所示。

\sourcefigure{fig09_06.pdf}{以人工标注的视频片段描述为线索，用多模态思维链检测情绪。第一步把多段描述统一用“他”作主语，去掉重复文字、合成完整句子并检查标点。第二步用示例引导情绪归纳，示例中“他看起来很开心，但实际上很焦虑”归为“焦虑”。}{fig:ch09-emotion-cot}

MM-CoT 的推理链配置大致有两类。
\begin{enumerate}
\item \textbf{自适应配置。}大语言模型自主确定推理链的长度（\refcite{ch09-ref-057}{Wu et al., 2023a}）。模型可根据任务的复杂程度调整推理过程，使响应更切合任务，也可能更准确。不同任务需要的推理深度相差很大时，这种方式尤其有用。
\item \textbf{预定义配置。}预先确定推理链长度（\refcite{ch09-ref-020}{Himakunthala et al., 2023}），在不同任务上用一致的推理结构。这样可以简化模型的运作，但在复杂情形中可能限制推理深度。
\end{enumerate}

除链长配置外，MM-CoT 还可选择不同的推理链生成方式。
\begin{enumerate}
\item \textbf{基于填补的方式。}推导能够在逻辑上衔接前后文的中间步骤，补上推理过程中的空缺（\refcite{ch09-ref-020}{Himakunthala et al., 2023}）。模型需找出并补齐思路中缺失的环节，使思路连贯且符合逻辑。例如，给定一组展示故事的图像，要求模型叙述事件，模型就要补全图像之间的情节，使故事连贯。
\item \textbf{基于预测的方式。}根据指令或此前的推理步骤等给定条件，将推理链向后延伸（\refcite{ch09-ref-057}{Wu et al., 2023a}）。模型需要理解当前上下文并推测合乎逻辑的后续内容，从而合成新的推理步骤。例如，要求模型预测图像故事的下一幕时，它需根据给定图像和文本描述延续情节，推测接下来会发生哪些事件或动作。
\end{enumerate}
表~\ref{tab:ch09-mm-cot}介绍若干常用的 MM-CoT 推理数据集。

\begin{table}[!htbp]
\centering\small
\caption{多模态思维链数据集。}
\label{tab:ch09-mm-cot}
\setlength{\tabcolsep}{3pt}
\begin{tabularx}{\linewidth}{@{}p{19mm}p{18mm}p{24mm}>{\raggedright\arraybackslash}X@{}}
\toprule
数据集 & 模态 & 规模 & 说明\\
\midrule
EMER & 视频—文本 & 100 & 侧重可解释的情绪推理，为推理任务提供线索和摘要。\\
EgoCOT & 视频—文本 & 3,670 小时 & 用于具身场景规划的大规模具身规划数据集。\\
VIP & 视频—文本 & 3.6M，测试集 1.5K & 用于视频思维链评估，提供大量描述数据和推理阶段的挑战任务。\\
ScienceQA & 图像—文本 & 21K 问答对 & 用于深入分析的多模态多项选择题数据集，覆盖科学及其他多个领域。\\
\bottomrule
\end{tabularx}
\end{table}

\subsection{多模态 RLHF}
多模态大语言模型需要整合和解释不同类型的数据，因此比在单一模态上训练的模型面临更多挑战。与其在单模态模型中的用法类似，RLHF 可以解决多模态大语言模型面临的多种问题，包括纳入人类偏好和选择、把人类反馈融入描述，以及生成符合安全与伦理标准的响应。下面介绍关于可信度以及如何纳入人类偏好和进行对齐的若干研究。

\begin{practicaltip}
Li et al. (2023) 研究如何通过偏好蒸馏，生成对用户有帮助且以视觉上下文为依据的响应。研究提出 VLFeedback 数据集，包含 80,000 条多模态指令、12 个大型视觉语言模型（LVLM）的响应，以及 GPT-4V 提供的偏好标注。用该数据集调优的 Silkie 在多种基准上大幅超过基础模型。与人工标注的数据集相比，该数据集有效提高了大型视觉语言模型的感知与认知能力，在扩展规模以及改善多项任务表现方面也有优势。
\end{practicaltip}

RLHF-V 是一种基于 RLHF 的方法，旨在通过细粒度的人类反馈对齐模型行为，提高多模态大语言模型的可信度（\refcite{ch09-ref-067}{Yu et al., 2023}）。它针对已有模型的一项重要问题，即模型容易生成缺乏关联图像事实依据的文本幻觉，影响在现实应用中的可靠性，在结果影响较大的使用场景中尤其如此。RLHF-V 通过纠正响应各片段中的幻觉收集人类偏好，并基于这些反馈进行密集的直接偏好优化。在五项基准上的自动评估和人工评估表明，RLHF-V 显著提高了多模态大语言模型行为的可信度，所需的数据和计算资源也较少。

\begin{practicaltip}
Sun et al. (2023) 提出“事实增强 RLHF”（Factually Augmented RLHF）对齐算法，将图像描述和多项选择题的正确答案等事实内容纳入已有奖励模型。这一做法旨在减少 RLHF 中的奖励投机，显著改善模型效果。研究还将已有的人工图文对纳入视觉指令调优的训练数据，该数据原先由 GPT-4 生成。这样可以改善模型的整体表现。该研究首次将 RLHF 用于语言多模态模型（LMM），使其在 LLaVA-Bench 上的表现明显提升，接近纯文本 GPT-4 的结果。
\end{practicaltip}

\subsection{输出投影器}
\label{sec:ch09-output-projector}
\begin{keypoint}
输出投影器 $\mathrm{OUT\_ALIGN}_{T\to X}$ 将大语言模型产生的信号词元表征 $S_X$ 转成后续模态生成器 $\mathrm{MG}_X$ 可以解释的特征 $H_X$。
\end{keypoint}

具体来说，给定模态与文本的配对数据集 $\{(I_X,t)\}$，大语言模型先处理输入 $t$，得到 $S_X$，随后将其转为 $H_X$。主要目标是使 $H_X$ 与模态生成器可理解的条件表征充分对齐。
\begin{equation}
\arg\min_{\mathrm{OUT\_ALIGN}_{T\to X}}L_{\mathrm{mse}}(H_X,\tau_X(t)),
\label{eq:ch09-output-projector-loss}
\end{equation}
其中
\begin{equation}
H_X=\mathrm{OUT\_ALIGN}_{T\to X}(S_X).
\label{eq:ch09-output-projector}
\end{equation}
$L_{\mathrm{mse}}$ 是均方误差损失，用于最小化投影特征 $H_X$ 与文本条件表征 $\tau_X(t)$ 的差异。$\tau_X$ 是 $\mathrm{MG}_X$ 中的文本条件编码器。此优化过程主要处理文本，无需直接提供音频或视觉等非文本模态 $X$ 的输入。

\begin{practicaltip}
输出投影器通常用小型 Transformer 或多层感知机（MLP）实现，以兼顾效率与适应性。原书在此处将 MLP 写作 MLLP。
\end{practicaltip}

\subsection{模态生成器}
\label{sec:ch09-modality-generator}
\begin{keypoint}
模态生成器 $\mathrm{MG}_X$ 用于产生各种模态的输出，将编码特征转为多模态内容。
\end{keypoint}

\begin{practicaltip}
该组件常用先进的潜空间扩散模型（latent diffusion model，LDM）合成图像、视频和音频等特定模态的输出（\refcite{ch09-ref-075}{Zhao et al., 2022}）。常用实现包括用于图像合成的 Stable Diffusion、用于视频合成的 Zeroscope 和用于音频生成的 AudioLDM-2（\refcite{ch09-ref-006}{Cerspense, 2023}; \refcite{ch09-ref-037}{Liu et al., 2023}; \refcite{ch09-ref-044}{Rombach et al., 2022}）。
\end{practicaltip}

用输出投影器给出的 $H_X$ 作为条件输入，引导生成高质量多模态内容所需的去噪步骤。训练阶段，先用预训练的变分自编码器（VAE）将原始内容编码为潜变量特征 $z_0$（\refcite{ch09-ref-026}{Kingma and Welling, 2013}）。然后加入噪声 $\epsilon$ 扰动潜在表征，得到带噪潜变量特征 $z_t$。

通常用预训练的 U-Net（$\epsilon_X$）计算条件 LDM 损失 $L_{X\text{-gen}}$（\refcite{ch09-ref-045}{Ronneberger et al., 2015}）。
\begin{equation}
L_{X\text{-gen}}:=
\mathbb{E}_{\epsilon\sim\mathcal{N}(0,1),\,t}
\left\lVert\epsilon-\epsilon_X(z_t,t,H_X)\right\rVert_2^2.
\label{eq:ch09-modality-generator-loss}
\end{equation}
通过最小化 $L_{X\text{-gen}}$ 优化 $\mathrm{IN\_ALIGN}_{X\to T}$ 和 $\mathrm{OUT\_ALIGN}_{T\to X}$。

% 原书第 9.4 节。
\section{评估基准}
\label{sec:ch09-benchmarks}
本节概述用于评估多模态大语言模型在不同模态与任务上表现的部分基准数据集，着重介绍任务种类较多、覆盖模态较广或在本领域广泛使用的基准。更详细的列表见 Yin et al. (2023)。

\begin{enumerate}
\item \textbf{CMMU} 收录 12,000 道多模态题目，由人工从大学考试、测验和教材中整理，涉及艺术与设计、商业、科学、健康与医学、人文与社会科学、技术与工程六个基础学科领域（\refcite{ch09-ref-071}{Zhang et al., 2024b}）。对应的 MMMU 则覆盖 30 个科目。CMMU 含有 39 种图像，包括统计图、示意图、地图、表格、乐谱和化学结构等，用于考察不同方面的多模态理解能力。

\item \textbf{MMCBench} 提供评估大型多模态模型（LMM）的详细框架，主要考察其在常见输入扰动下的稳健性和自洽性（\refcite{ch09-ref-073}{Zhang et al., 2024c}）。它关注文本、图像和语音模态的相互作用，涉及文本生成图像、图像生成文本、文本生成语音和语音转文本等主要任务。

\item \textbf{MMVP} 通过视觉问答任务评估多模态大语言模型的视觉能力（\refcite{ch09-ref-051}{Tong et al., 2024}）。它含有一个包含 300 张测试图像的目录，以及一个包含问题和正确答案的 CSV 文件。

\item \textbf{TimeIT} 涉及六种与时间戳相关的视频任务，纳入 12 个来自不同领域的数据集（\refcite{ch09-ref-043}{Ren et al., 2023}）。它侧重时间敏感的长视频理解任务，包括密集视频描述、视频时刻定位、视频摘要、精彩片段检测、步骤定位和语音转录，训练数据共 124,861 个实例。

\item \textbf{ViP-Bench} 包含 303 个来自 MM-Vet、MMBench 和 Visual Genome 的图像与问题配对，旨在测试多模态模型的视觉推理能力（\refcite{ch09-ref-004}{Cai et al., 2023}）。它在图像区域层面考察视觉理解的六个方面，即识别、光学字符识别（OCR）、知识、数学、物体关系推理和语言生成。基准用 GPT-4 对模型响应给出 0 至 10 分，便于定量比较。

\item \textbf{M3DBench} 汇编了超过 320K 对三维多模态指令遵循数据，包括具有不同多模态提示的 138K 多条指令（\refcite{ch09-ref-004}{Cai et al., 2023}）。它利用已有数据集以及由大语言模型为各种三维任务生成的指令。覆盖从物体检测到问答的任务，每类任务有专门的指令和响应。它通过模式匹配过滤不相关的响应，以改善数据质量，供三维指令遵循评估使用。

\item \textbf{Video-Bench} 提供评估视频大语言模型的综合基准（\refcite{ch09-ref-041}{Ning et al., 2023}）。它包含十项任务，用于考察对视频内容的理解、借助已有知识回答问题以及理解和决策等能力，共约 15,033 个实例。

\item \textbf{Bingo} 对多模态理解中模型失败与成功的实例进行分类，包括 190 个失败实例和 131 个成功实例（\refcite{ch09-ref-013}{Cui et al., 2023}）。每个实例配有一两个问题，归入“干扰”（图像对图像、文本对图像）或“偏见”（区域、OCR、事实）类别。该基准旨在分析响应出现幻觉的具体原因，用文化、语言与事实背景各异的图像，详细考察 GPT-4V(ision) 中的偏见。

\item \textbf{MMHAL-BENCH} 使用专门的指标评估多模态模型的幻觉（\refcite{ch09-ref-049}{Sun et al., 2023}）。它包含 96 个图像与问题配对，涉及八类问题和十二种物体主题，用来考察大型多模态模型是否会对图像内容作出不实陈述。为避免数据泄漏，基准采用 Open Images 验证集和测试集中的图像，覆盖“配件”“动物”“交通工具”等物体大类。GPT-4 比较模型的输出、人工编写的答案与图像内容，判断是否出现幻觉。

\item \textbf{Sparkles} 利用 GPT-4 构建多模态对话数据集，模拟围绕图像和文本展开的真实对话，涵盖多种现实场景（\refcite{ch09-ref-024}{Huang et al., 2023}）。生成过程采用两轮对话模式，先由用户就图像提问，再由助手详细回答，然后引入新图像供后续讨论。数据生成时用对话示例进行上下文学习，并根据候选图像的描述选择相关图像。当时的 GPT-4 仅接受文本输入，因此用详细的文本描述代表图像。

\item \textbf{SciGraphQA} 是一个围绕学术图表构建的大规模合成多轮问答数据集（\refcite{ch09-ref-033}{Li and Tajbakhsh, 2023}）。它含有 295,000 个样本，源自 arXiv 上 2010 至 2020 年的 290,000 篇计算机科学或机器学习论文。用 PaLM-2 根据论文中的图表生成对话，并纳入标题、摘要、相关段落及上下文数据。每段对话平均有 2.23 轮问答。在 3,000 个样本的测试集上，GPT-4 给出的问答匹配质量平均为 8.7 分，满分为 10 分。

\item \textbf{LAMM} 是一个多模态指令调优数据集，用多种视觉任务中的图像与点云，构建 186,000 对语言与图像、10,000 对语言与三维数据的指令和响应（\refcite{ch09-ref-064}{Yin et al., 2024}）。数据集利用 GPT API 和自指令方法构建，包括日常对话、事实知识对话、图像与三维场景的详细描述、视觉任务对话四类配对，以增强视觉任务泛化能力。它涉及多种二维和三维视觉任务，例如描述生成、场景图识别、视觉问答、分类、检测、计数和 OCR。
\end{enumerate}

% 原书第 9.5 节。
\section{代表性的多模态大语言模型}
\label{sec:ch09-models}
本节概述若干当时先进的多模态大语言模型，展示它们如何把各种模态整合进框架。表~\ref{tab:ch09-models}列出常用模型，将各自的组件对应到前文的通用框架。接下来深入讨论三个模型，它们代表了处理多模态数据的复杂程度和能力的进展。首先讨论结合视觉与语言的 Flamingo，再讨论可处理文本以及视频和音频的 Video-LLaMA，最后讨论能够在多种模态之间转换的 NExT-GPT。\footnote{原书此处称 NExT-GPT 支持六种模态。其编码器 ImageBind 可处理六种模态，而 \href{https://arxiv.org/abs/2309.05519}{NExT-GPT 论文}描述的端到端输入与输出为文本、图像、视频和音频四种，与表~\ref{tab:ch09-models}和第~\ref{sec:ch09-next-gpt}~节相符。}

% 横向版面保留原书七列与 18 个模型。
\begin{landscape}
\small\setlength{\tabcolsep}{3pt}\renewcommand{\arraystretch}{1.09}
\begin{longtable}{@{}>{\raggedright\arraybackslash}p{21mm}
 >{\raggedright\arraybackslash}p{35mm}
 >{\raggedright\arraybackslash}p{38mm}
 >{\raggedright\arraybackslash}p{26mm}
 >{\raggedright\arraybackslash}p{27mm}
 >{\raggedright\arraybackslash}p{14mm}
 >{\raggedright\arraybackslash}p{29mm}@{}}
\caption{多模态大语言模型的输入与输出模态及其组件。I 为图像，V 为视频，A 为音频，T 为文本，3D 为三维数据，PT 为预训练，IT 为指令调优。}\label{tab:ch09-models}\\
\toprule
模型 & 输入到输出 & 模态编码器 & 输入投影器 & 核心大语言模型 & 输出投影器 & 模态生成器\\
\midrule
\endfirsthead
\multicolumn{7}{@{}l}{表 \thetable\ （续）}\\
\toprule
模型 & 输入到输出 & 模态编码器 & 输入投影器 & 核心大语言模型 & 输出投影器 & 模态生成器\\
\midrule
\endhead
\midrule
\multicolumn{7}{r@{}}{下页续}\\
\endfoot
\bottomrule
\endlastfoot
Flamingo & I+V+T$\to$T & I/V：NFNet-F6 & 交叉注意力 & Chinchilla-\allowbreak 1.4B/7B/70B & — & —\\
BLIP-2 & I+T$\to$T & I：CLIP/EVA-CLIP ViT@224 & Q-Former 与线性投影器 & Flan-T5/OPT & — & —\\
LLaVA & I+T$\to$T & I：CLIP ViT-L/14 & 线性投影器 & Vicuna-7B/13B & — & —\\
IDEFICS & I+T$\to$T & I：OpenCLIP ViT-H/14 & 交叉注意力 & LLaMA-v1 7B/65B & — & —\\
MiniGPT-4 & I+T$\to$T & I：EVA-CLIP ViT-G/14 & Q-Former 与线性投影器 & Vicuna-13B & — & —\\
X-LLM & I+V+A+T$\to$T & I/V：ViT-G，A：C-Former & Q-Former 与线性投影器 & ChatGLM-6B & — & —\\
VideoChat & V+T$\to$T & V：ViT-G & Q-Former 与线性投影器 & Vicuna & — & —\\
InstructBLIP & I+V+T$\to$T & I/V：ViT-G/14@224 & Q-Former 与线性投影器 & Flan-T5/Vicuna & — & —\\
Video-LLaMA & I+V+A+T$\to$T & I/V：EVA-CLIP ViT-G/14，A：ImageBind & Q-Former 与线性投影器 & Vicuna/LLaMA & — & —\\
BuboGPT & I+A+T$\to$T & I：CLIP/EVA-CLIP ViT，A：ImageBind & Q-Former 与线性投影器 & Vicuna（冻结） & — & —\\
Qwen-VL-\allowbreak (Chat) & I+T$\to$T & I：以 OpenCLIP ViT-bigG 初始化的 ViT@448 & 交叉注意力 & Qwen-7B（PT 冻结，IT 用 PEFT） & — & —\\
PaLM-E & I+3D+T$\to$T & I：ViT，3D：OSRT & 仿射变换 & PaLM（PT 冻结并协同训练） & — & —\\
MACAW-LLM & I+V+A+T$\to$T & I/V：CLIP，A：Whisper & 线性投影器 & Llama-7B & — & —\\
NExT-GPT & I+V+A+T$\to$\allowbreak I+V+A+T & I/V/A：ImageBind & 线性投影器 & Vicuna-7B（PEFT） & 小型 Trans\-former & I：Stable Diffusion，V：Zeroscope，A：AudioLDM\\
MiniGPT-5 & I+T$\to$I+T & I：EVA-CLIP ViT-G/14 & Q-Former 与线性投影器 & Vicuna-7B（PEFT） & 小型 Trans\-former 与 MLP & I：Stable Diffusion 2\\
LLaVA-1.5 & I+T$\to$T & I：CLIP ViT-L@336 & MLP & Vicuna-v1.5-\allowbreak 7B/13B（PT 冻结，IT 用 PEFT） & — & —\\
X-InstructBLIP & I+V+A+3D+T$\to$T & I/V：EVA-CLIP ViT-G/14，A：BEATs，3D：ULIP-2 & Q-Former 与线性投影器 & Vicuna-v1.1-\allowbreak 7B/13B（冻结） & — & —\\
CoDi-2 & I+V+A+T$\to$\allowbreak I+V+A+T & I/V/A：ImageBind & MLP & Llama-2-Chat-7B（PT 冻结，IT 用 PEFT） & MLP & I：Stable Diffusion 2.1，V：Zeroscope-v2，A：AudioLDM-2\\
\end{longtable}
\end{landscape}


\subsection{Flamingo（图像、视频与文本）}
Flamingo 是早期的视觉语言模型，在多项视觉与语言任务上推进了少样本学习（\refcite{ch09-ref-001}{Alayrac et al., 2022}）。在 16 项任务中，它有 6 项超过经过微调的先进模型，同时使用的任务专用训练数据少得多。下面按通用框架的组件介绍它的架构。该模型没有使用多模态 RLHF，输出为文本，因此也无需输出投影器和其他模态的生成器。

\sourcefigure{fig09_07.pdf}{Flamingo 的架构。冻结的视觉编码器提取图像和视频特征，由 Perceiver 重采样器转换成定长视觉词元。可训练的门控交叉注意力层与冻结的语言模型层交错排列，以视觉特征为条件生成文本。}{fig:ch09-flamingo}

\subsubsection{模态编码器}
Flamingo 的视觉编码器是 F6 规格的无归一化 ResNet（Normalizer-Free ResNet，NFNet），通过视觉与文本的对比学习来编码视觉输入。在对比预训练阶段，底本描述其文本编码器为 BERT，有别于使用 GPT-2 的做法。对视觉和文本的嵌入进行均值池化，再投影到联合嵌入空间，便于整合两种模态。Flamingo 生成文本时使用的核心语言模型见下文。

\subsubsection{输入投影器}
Flamingo 可处理图像和视频，其视觉特征的数量也随输入而变。Perceiver 重采样器（Perceiver Resampler）把这些特征统一转为 64 个视觉词元，见图~\ref{fig:ch09-flamingo}。将门控交叉注意力层（GATED XATTN-DENSE）插入冻结的语言模型层之间，可实现语言与视觉模态的对齐，使模型在生成文本词元时关注视觉词元。

\subsubsection{预训练：核心大语言模型、数据集和任务目标}
Flamingo 以 Chinchilla 为基础，冻结预训练的语言模型层。\footnote{原书此处写作“冻结九个预训练的 Chinchilla 层”。\href{https://arxiv.org/abs/2204.14198}{Flamingo 论文}第 2.2 节说明冻结预训练语言模型，训练新插入的交叉注意力层。}训练使用四个不同的数据集，即 M3W（交错图文）、ALIGN（图文对）、LTIP（图文对）和 VTP（视频与文本对）。Flamingo 同时根据前序文本和视觉词元预测后续文本词元 $y$，其概率为
\begin{equation}
p(y\mid x)=\prod_{\ell=1}^{L}p(y_\ell\mid y_{<\ell},x_{\leq\ell}).
\label{eq:ch09-flamingo-probability}
\end{equation}
训练损失是各数据集生成文本的期望负对数似然之加权和。
\begin{equation}
\sum_{m=1}^{M}\lambda_m\,
\mathbb{E}_{(x,y)\sim D_m}\left[
-\sum_{\ell=1}^{L}\log p(y_\ell\mid y_{<\ell},x_{\leq\ell})
\right],
\label{eq:ch09-flamingo-loss}
\end{equation}
其中 $D_m$ 是第 $m$ 个数据集，$\lambda_m$ 是相应的训练权重。

\subsubsection{多模态大语言模型的调优与增强}
Flamingo 在上下文学习中表现出色。它使用同一套模型权重和仅 32 个任务专用示例，少样本表现就超过了针对具体任务微调的先进模型。其示例数比当时先进方法所用的任务专用训练样本少约三个数量级。分析中把图像或视频（视觉输入）和对应文本（预期响应或问题等任务相关信息）配成示例，用来预测新视觉查询的响应。非问答任务的默认提示是 \texttt{Output: output}，问答或视觉对话任务的默认提示是 \texttt{Question: question Answer: answer}。

\subsection{Video-LLaMA（图像、视频、音频与文本）}
Zhang et al. (2023) 提出 Video-LLaMA 多模态框架，使大语言模型能理解视频中的视觉和听觉要素。此前的方法已让大语言模型处理视觉或音频信号。Video-LLaMA 利用冻结的预训练视觉编码器和音频编码器，以及冻结的大语言模型，开展跨模态训练，综合处理视听信息。它侧重视频理解，处理两个问题，即捕捉视觉场景随时间的变化，以及融合音视频信息。实验显示，它能基于音频和视频进行对话，可作为视听人工智能助手的原型。下面按通用框架的组件讨论其架构。模型也未使用多模态 RLHF，输出仅为文本，因此无需输出投影器和其他模态的生成器。

\subsubsection{模态编码器}
视觉分支采用冻结的视觉编码器，其中包括 EVA-CLIP 的 ViT-G/14 和 BLIP-2 的 Q-Former，以处理视频帧，见图~\ref{fig:ch09-video-llama}。每帧转换成一组图像嵌入向量，得到帧表征序列 $V=[v_1,v_2,\ldots,v_N]$。其中 $v_i\in\mathbb{R}^{K_f\times d_f}$ 表示第 $i$ 帧的 $K_f$ 个 $d_f$ 维图像嵌入。

视频的音频部分采用预训练的 ImageBind 作为音频编码器（\refcite{ch09-ref-018}{Girdhar et al., 2023}）。从视频中均匀采样 $M$ 个长度为 2 秒的音频片段。将各片段转成带 128 个梅尔频带的频谱图，以提取频谱特征。音频编码器再把各频谱图转成稠密向量表征。因此，整个视频的音频表征为 $A=[a_1,a_2,\ldots,a_M]$，其中 $a_i$ 是第 $i$ 个音频片段经编码后的特征向量。

\sourcefigure{fig09_08.pdf}{Video-LLaMA 的架构。视觉与音频分支分别编码图像或视频帧以及音频片段。两条分支各用位置嵌入、Q-Former 和线性层，将输入转成大语言模型可处理的软提示。}{fig:ch09-video-llama}

\subsubsection{输入投影器}
在视频和音频分支中，都使用 Q-Former 与线性投影相结合的方式，将编码后的模态与文本数据对齐。

在视觉语言分支中，帧表征 $v_i$ 来自冻结的图像编码器，本身不带时间信息。因此先添加位置嵌入，再将加了位置编码的帧表征输入视频 Q-Former。它融合各帧表征，产生 $k_V$ 个 $d_v$ 维的视频嵌入向量，组成视频表征 $\hat{v}\in\mathbb{R}^{k_V\times d_v}$，同时捕捉视觉内容和时间变化。然后用线性层把视频嵌入向量转换为视频查询向量，与大语言模型的文本嵌入具有相同维度。前向传播时，将视频查询向量与文本嵌入拼接起来，作为视频软提示。这样，冻结的大语言模型可以在视频内容的条件下生成文本输出，将视频信息纳入多模态理解过程。

音频语言分支采用类似方法。位置嵌入层为音频片段添加时间信息，模型因而能捕捉理解视频中声音的顺序和变化所需的时间特征。随后，音频 Q-Former 将不同片段的特征合成统一的音频表征，再由线性层把它映射到大语言模型的嵌入空间。

\subsubsection{预训练：核心大语言模型、数据集和任务目标}
Video-LLaMA 以 Vicuna-7B 为核心大语言模型，实现多模态理解和生成。

Video-LLaMA 预训练的视觉语言分支使用 WebVid-2M，其中包含来自素材网站的短视频及其文本描述。另一个数据集是从 CC3M 派生、经 Liu et al. (2024) 筛选的 CC595k 图像描述数据集。这两个数据集为预训练中的视频到文本生成任务提供数据。音频语言分支的 ImageBind 音频编码器已经在多种模态之间实现对齐，因此不需要这一步预训练。

\subsubsection{多模态大语言模型的调优与增强}
完成视频语言分支预训练后，Video-LLaMA 已能较好地根据视频信息生成内容，但遵循具体指令的能力还有待增强。因此，研究人员进行了指令微调，使用以下数据集。
\begin{enumerate}
\item 来自 LLaVA 数据集的 150K 条基于图像的指令。
\item 来自 MiniGPT-4 的 3K 条基于图像的指令。
\item 来自 VideoChat 的 11K 条基于视频的指令。
\end{enumerate}
调优音频分支时，为了克服音频与文本配对数据不足的问题，研究人员也把上述视觉语言数据集纳入训练，使相关组件能够学习如何将 ImageBind 编码器产生的共享嵌入空间对齐到大语言模型的嵌入空间。

\subsection{NExT-GPT（任意模态组合到任意模态组合）}
\label{sec:ch09-next-gpt}
NExT-GPT 是一种通用多模态大语言模型，将大语言模型与多模态适配器、扩散解码器结合，可以处理并生成文本、图像、视频和音频（\refcite{ch09-ref-059}{Wu et al., 2023c}）。微调只更新少量参数，训练成本较低，也便于扩展到新的模态。系统还采用模态切换指令调优及相应数据集，改善跨模态理解和生成，表明可以构建适用于多种模态的统一人工智能智能体。

\subsubsection{模态编码器}
NExT-GPT 使用 ImageBind 作为各输入模态的通用编码器。以往许多研究为不同模态分别采用专用编码器。ImageBind 可以建立涵盖多种模态的联合嵌入空间，无需对每一种可能的模态组合都提供训练数据，也能达到当时先进的表现。

\subsubsection{输入投影器}
NExT-GPT 使用一个有 400 万参数的线性投影层，把经 ImageBind 得到的输出转化为类似语言的表征，将各模态对齐到大语言模型易于理解和处理的格式。

\subsubsection{预训练：核心大语言模型、数据集和任务目标}
NExT-GPT 的核心大语言模型是 Vicuna-7B，一种常用于多模态大语言模型的开源文本大语言模型。\footnote{底本正文两处写作 Vicuna2。这里按表~\ref{tab:ch09-models}及 \href{https://arxiv.org/abs/2309.05519}{NExT-GPT 论文}的模型名称写作 Vicuna-7B。}

预训练第一阶段冻结流水线中处理输入的编码器、处理输出的解码器以及大语言模型，只通过反向传播训练用于输入对齐的线性投影层。目标是用各模态专用的数据集，把图像、音频或视频与相应的文本描述对齐。图像模态对齐使用 CC3M，包含 300 万个图像与描述配对。视频模态对齐使用 WebVid-10M，包含 1,070 万个视频与描述配对，内容来自素材网站，总计 52,000 小时。音频模态使用 AudioCaps，其中约有 46,000 个音频片段和人工编写文本的配对，由众包方式在 AudioSet 数据集上整理。通过这种方法，NExT-GPT 学会根据输入模态生成描述，并与人工给定的参考描述匹配。

\begin{figure}[!htbp]
\centering
\includegraphics[trim=0 173bp 0 172bp,clip,width=\linewidth]{assets/fig09_09.pdf}
\caption{NExT-GPT 的架构。输入端用 ImageBind 和线性投影编码非文本数据。输出端把大语言模型产生的模态信号词元经相应的输出投影器转换，再交给图像、视频或音频扩散解码器。}
\label{fig:ch09-next-gpt}
\end{figure}

训练第二阶段重点训练输出投影层。此时冻结 ImageBind、大语言模型及输入投影层，并继续使用第一阶段的数据集，即图像用 CC3M，视频用 WebVid-10M，音频用 AudioCaps。

\subsubsection{多模态大语言模型的调优与增强}
NExT-GPT 训练的第三阶段采用多模态指令调优，以改善模型理解和执行涉及不同模态的复杂指令的能力。此阶段以 LoRA 微调核心大语言模型 Vicuna-7B，同时微调输入和输出投影层，编码器和解码器仍保持冻结。所用数据集如下。
\begin{enumerate}
\item \textbf{“文本与 X 到文本”数据。}这里的 X 指 NExT-GPT 支持的任一非文本模态，即图像、视频或音频。将文本与其中一种模态组合后送入模型，模型生成能准确反映组合输入的文本响应。采用的数据集包括视觉语言数据集 LLaVA-Instruct-150K、MiniGPT-4 的图像描述数据集及 VideoChat 的视频指令数据集。
\item \textbf{“文本到文本与 X”数据。}这一数据集用于训练模型根据文本生成文本及其他模态的内容。研究人员从 Conceptual Captions、MS COCO（Microsoft Common Objects in Context）、AudioCaps 等已有语料和基准收集“X 与描述”配对，构建文本到多模态（T2M）数据。再用模板和 GPT-4 编写包含这些描述的多种文本指令，得到可用于训练模型从文本提示同时生成文本和多模态内容的数据集。
\item \textbf{MosIT 数据。}NExT-GPT 的一项创新是开发模态切换指令调优数据集（Modality-switching Instruction Tuning，MosIT），以提高模型跨模态遵循指令的能力。已有数据集未能充分反映用户与人工智能通过不同媒介交互时的复杂性，因此需要构造更丰富的数据。研究团队从外部来源收集材料，包括从 YouTube 收集视频，并用 Stable-XL 和 Midjourney 等人工智能生成内容（AIGC）工具制作图像和音频片段。MosIT 的每段对话有 3 至 7 轮，在人机交互的输入与输出之间切换模态，最终得到 5,000 段对话。
\end{enumerate}

\subsubsection{输出投影器}
NExT-GPT 的输出投影器把大语言模型生成的词元转换为各模态解码器可处理的格式。为此，NExT-GPT 采用小型 Transformer（3100 万参数），针对各模态分别进行转换。这些投影器在整体训练过程的第二、第三阶段训练。

\subsubsection{模态生成器}
NExT-GPT 最后用专用解码器产生不同模态的输出。系统收到大语言模型生成的多模态信号和指令，经基于 Transformer 的层转换成解码器可以处理的格式。NExT-GPT 为各模态采用先进的扩散模型，图像用 Stable Diffusion，视频用 Zeroscope，音频用 AudioLDM。这些模型作为条件扩散模型集成到系统中，用转换后的信号表征作条件，生成所指定模态的内容。

% 原书第 9.6 节。
\section{教程：微调图像到文本的多模态大语言模型}
\label{sec:ch09-tutorial}

\subsection{概述}
讨论过多模态大语言模型的理论基础后，可以测试“文本与 X 到文本”的模型。本例用图像作为模态 X。图像与文本到文本的模型可用于识别图像特征、对图像中的事件分类和自动生成描述等任务。\footnote{原书此处误写为图像与文本到图像。按本节所用模型及任务，输出模态为文本。}本教程测试一种当时先进的多模态大语言模型在图像分类和描述生成方面的开箱能力，再考察如何通过微调与少样本提示改善表现。

我们在同一个数据集上安排两个实验。首先，要求模型识别每张图像中的运动项目，比较零样本和微调两种方案的结果。然后让模型用简短文字描述图像中发生的事，比较零样本、少样本和微调三种方式。

\begin{keypoint}
\textbf{学习目标}
\begin{itemize}
\item 加载 90 亿参数的 IDEFICS 模型，用任意文本和图像组成的提示调用模型。
\item 在 100SIC 测试集上生成零样本预测，尝试通过 QLoRA 微调改善表现。
\item 为 100SIC 测试集生成零样本描述，并与微调和上下文学习生成的描述比较。
\end{itemize}
\end{keypoint}

本节为教程的简版，完整版见
\url{https://github.com/springer-llms-deep-dive/llms-deep-dive-tutorials}。

\subsection{实验设计}
可用的多模态大语言模型很多。为了缩小范围，我们考察能够在 Google Colab 笔记本中用 QLoRA 调优、且已与 Hugging Face 集成的模型，以便直接使用其 PEFT 和微调流程。基于这些条件，我们选择 90 亿参数的 IDEFICS（Image-aware Decoder Enhanced à la Flamingo with Interleaved Cross-attentionS）。它仿照 Flamingo 构建，是一种将文本和图像作为输入、输出文本的大语言模型（\refcite{ch09-ref-028}{Laurençon et al., 2023}）。模型可接受任意交错排列的文本和图像，并输出文本响应。

实验采用 Kaggle 上的 100 项运动图像分类数据集（100 Sports Image Classification，100SIC）。\footnote{\url{https://www.kaggle.com/datasets/gpiosenka/sports-classification/data}}该数据集包含许多标有运动项目的小尺寸照片，共涉及 100 项运动，包含约 13,000 张训练图像，以及各 500 张测试和验证图像。为微调描述生成任务，我们补充使用 Flickr30k 数据集（\refcite{ch09-ref-066}{Young et al., 2014}）的一个子集。Flickr30k 收录了 30,000 多对图像与描述。这里采用 Shin Thant 筛出的 Flickr30k 运动图像子集。\footnote{\url{https://github.com/ShinThant3010/Captioning-on-Sport-Images}}

\subsection{结果与分析}
\subsubsection{预测运动项目}
首先加载模型。IDEFICS 对于单块中等配置的 GPU 来说过大，难以高效地预测和调优。因此用 BitsAndBytes 将模型量化到 4 位，并采用 QLoRA 微调。运动分类采用以下提示模板。

\noindent\begin{minipage}{\linewidth}
\begin{lstlisting}[caption={运动项目分类提示。},label={lst:ch09-sport-prompt}]
<image>
Question: What sport is in this image?
Answer:
\end{lstlisting}
\end{minipage}\par

用这个模板预测测试集中的每张图像，再将输出与数据集编制者给出的标签比较。

\noindent\begin{minipage}{\linewidth}
\begin{lstlisting}[caption={测试集零样本预测结果。},label={lst:ch09-sport-zero}]
- Zero-shot results:
  - 212 / 500 correct
\end{lstlisting}
\end{minipage}\par

模型在约 42\% 的图像上猜对了运动项目的名称。这里采用简单的精确匹配评估，所以即便模型猜出同一运动项目的另一个正确名称，也会记作预测错误。可在训练集上微调模型，改善预测。Llama 基础模型太大，不适合全参数微调，因此采用与第~4.6~节教程类似的 QLoRA 方法。每项运动选 10 个训练样本，用零样本示例中的模板构造问答对，微调模型，再对测试集进行预测。

\noindent\begin{minipage}{\linewidth}
\begin{lstlisting}[caption={微调后的测试集预测结果。},label={lst:ch09-sport-tuned}]
- Fine-tune results:
  - 419 / 500 correct
\end{lstlisting}
\end{minipage}\par

正确率从约 42\% 提高到约 84\%。表~\ref{tab:ch09-sport-outputs}选取了若干例子，具体说明这次改进的表现。

\begin{table}[!htbp]
\centering\fontsize{8.8pt}{11.5pt}\selectfont
\caption{零样本和微调在三个特选运动类别上的预测。雪车项目（bobsled）中，零样本模型多能识别运动，却常用不同名称。马车赛（chuckwagon racing）较少见，模型把它猜成其他马术比赛。拔河（tug of war）中，微调使这组预测变差，继续微调可能改善结果。ZS 表示零样本，FT 表示微调。}
\label{tab:ch09-sport-outputs}
\setlength{\tabcolsep}{2.1pt}
\renewcommand{\arraystretch}{1.14}
\begin{tabular}{@{}r*{6}{>{\raggedright\arraybackslash}p{16.6mm}}@{}}
\toprule
& \multicolumn{2}{c}{bobsled} & \multicolumn{2}{c}{chuckwagon racing} & \multicolumn{2}{c}{tug of war}\\
\cmidrule(lr){2-3}\cmidrule(lr){4-5}\cmidrule(lr){6-7}
序号 & ZS & FT & ZS & FT & ZS & FT\\
\midrule
1 & bobsledding & bobsled & rodeo & chuckwagon racing & hurling & oxen pulling\\
2 & the u & bobsled & horse-drawn carriage racing & chuckwagon racing & rugby & tug of war\\
3 & bobsleigh & bobsled & calgary stampede rodeo & chuckwagon racing & tug of war & log rolling\\
4 & bobsledding & bobsled & horseback riding & chuckwagon racing & tug of war & log rolling\\
5 & bobsled & bobsled & chariot racing & chuckwagon racing & tug of war & axe throwing\\
\bottomrule
\end{tabular}
\end{table}

\subsubsection{为照片生成描述}
图像到文本模型的另一种常见用途是自动生成图像描述。这里测试 IDEFICS 完成该任务的能力。与前面一样，可用一个简单的零样本提示模板要求模型生成描述。本练习选用一张美式橄榄球比赛的图像，其中身穿黑色球衣的维克森林大学四分卫正在传球。提示模板如下。

\noindent\begin{minipage}{\linewidth}
\begin{lstlisting}[caption={运动照片描述生成提示。},label={lst:ch09-caption-prompt}]
<image>
Question: What is a caption for this photo?
Answer:
\end{lstlisting}
\end{minipage}\par

使用零样本提示时，得到以下响应。

\noindent\begin{minipage}{\linewidth}
\begin{lstlisting}[caption={运动照片描述的零样本输出。},label={lst:ch09-caption-zero}]
Question: What is a caption for this photo? Answer: Aaron Murray,
     Georgia Bulldogs quarterback, throws a pass during the first
     half of the Chick-fil-A Bowl NCAA college football game
     against the Nebraska Cornhuskers
\end{lstlisting}
\end{minipage}\par

照片确实显示一名四分卫正在传球，但响应中的其他细节都是错的。球员既不是阿龙·默里，也不属于佐治亚斗牛犬队。比赛也不是 Chick-fil-A 碗赛，对手并非内布拉斯加玉米剥壳人队。这些信息都是模型的幻觉。尤其是最后两项错误事实，仅凭图像本身就无法确定。理想的描述应直接说明图像中可以看见什么，少写照片上无法清楚辨认的身份信息。

用上下文示例改善描述的成本较低。采用这种方法时，在目标图像之前提供几张训练图像及其人工编写的描述。得到的输出如下。

\noindent\begin{minipage}{\linewidth}
\begin{lstlisting}[caption={运动照片描述的上下文提示。},label={lst:ch09-caption-icl}]
<image1> Question: What is a caption for this photo? Answer: A
    man prepares to throw an ax at a target.
<image2> Question: What is a caption for this photo? Answer: A
    woman rolls a bowling ball down a bowling alley.
...
<image5> Question: What is a caption for this photo? Answer:
A man in a white jersey throws a football.
\end{lstlisting}
\end{minipage}\par

少样本条件下，模型生成了“一名穿白色球衣的男子掷出橄榄球”。球衣实际为黑色，因此仍有一处小错误。不过，模型采用了我们希望的描述形式，也没有凭空补充球员身份或对手之类的细节。这个办法有改进空间。

更昂贵的做法是用 Flickr30k 的运动图像和描述配对子集来微调。我们沿用上述 QLoRA 方法，采用零样本示例中的模板，用约 1,600 个样本微调 IDEFICS 基础模型。训练完成后，再为测试图像生成描述。

\noindent\begin{minipage}{\linewidth}
\begin{lstlisting}[caption={微调后的运动照片描述输出。},label={lst:ch09-caption-tuned}]
A football player in a black uniform is throwing a football.
\end{lstlisting}
\end{minipage}\par

这次的输出与少样本输出一样简洁，也符合照片内容，即“一名身穿黑色队服的橄榄球运动员正在掷球”。最后，用三种方法分别为 20 张测试图像生成描述，并从准确性和风格两方面人工评定。以下是结果。

\noindent\begin{minipage}{\linewidth}
\begin{lstlisting}[caption={测试集图像描述生成结果。},label={lst:ch09-caption-results}]
- Zero-shot results:
  - 7 / 20 acceptable
- In-context results:
  - 11 / 20 acceptable
- Fine-tune results:
  - 14 / 20 acceptable
\end{lstlisting}
\end{minipage}\par

\subsection{结论}
中等规模、以文本和图像为输入并输出文本的多模态大语言模型，已经具备相当的零样本能力，针对任务的微调可以大幅改善其表现。本教程说明，即使只做少量优化并使用常用参数，采用参数高效微调也能显著改善图像分类和开放式描述生成。生产应用需要更仔细地选择调优参数、训练集特性和多模态大语言模型架构。本练习表明，用现有的开源软件、投入较少工作，也可以构建一个表现尚可的图像分类器。

% PDF 441--447。
\begin{chapterreferences}
\refentry{ch09-ref-001}{Jean-Baptiste Alayrac, Jeff Donahue, Pauline Luc, Antoine Miech, Iain Barr, Yana Hasson, Karel Lenc, Arthur Mensch, Katherine Millican, Malcolm Reynolds, et al. Flamingo: a visual language model for few-shot learning. \textit{Advances in Neural Information Processing Systems}, 35:23716–23736, 2022.}
\refentry{ch09-ref-002}{Lalit Bahl, Peter Brown, Peter De Souza, and Robert Mercer. Maximum mutual information estimation of hidden markov model parameters for speech recognition. In \textit{ICASSP’86. IEEE international conference on acoustics, speech, and signal} \textit{processing}, volume 11, pages 49–52. IEEE, 1986.}
\refentry{ch09-ref-003}{Jinze Bai, Shuai Bai, Yunfei Chu, Zeyu Cui, Kai Dang, Xiaodong Deng, Yang Fan, Wenbin Ge, Yu Han, Fei Huang, et al. Qwen technical report. \textit{arXiv preprint} \textit{arXiv:2309.16609}, 2023.}
\refentry{ch09-ref-004}{Mu Cai, Haotian Liu, Siva Karthik Mustikovela, Gregory P Meyer, Yuning Chai, Dennis Park, and Yong Jae Lee. Making large multimodal models understand arbitrary visual prompts. \textit{arXiv preprint arXiv:2312.00784}, 2023.}
\refentry{ch09-ref-005}{Jean Carletta, Simone Ashby, Sebastien Bourban, Mike Flynn, Mael Guillemot, Thomas Hain, Jaroslav Kadlec, Vasilis Karaiskos, Wessel Kraaij, Melissa Kronenthal, et al. The ami meeting corpus: A pre-announcement. In \textit{International} \textit{workshop on machine learning for multimodal interaction}, pages 28–39. Springer, 2005.}
\refentry{ch09-ref-006}{Cerspense. Zeroscope: Diffusion-based text-to-video synthesis, 2023.}
\refentry{ch09-ref-007}{Feilong Chen, Minglun Han, Haozhi Zhao, Qingyang Zhang, Jing Shi, Shuang Xu, and Bo Xu. X-llm: Bootstrapping advanced large language models by treating multi-modalities as foreign languages. \textit{arXiv preprint arXiv:2305.04160}, 2023.}
\refentry{ch09-ref-008}{Sanyuan Chen, Yu Wu, Chengyi Wang, Shujie Liu, Daniel Tompkins, Zhuo Chen, and Furu Wei. Beats: Audio pre-training with acoustic tokenizers. \textit{arXiv preprint} \textit{arXiv:2212.09058}, 2022.}
\refentry{ch09-ref-009}{Wei-Lin Chiang, Zhuohan Li, Zi Lin, Ying Sheng, Zhanghao Wu, Hao Zhang, Lianmin Zheng, Siyuan Zhuang, Yonghao Zhuang, Joseph E Gonzalez, et al. Vicuna: An open-source chatbot impressing gpt-4 with 90\%* chatgpt quality. See \url{https://vicuna.lmsys.org} (accessed 14 April 2023), 2023.}
\refentry{ch09-ref-010}{Huyen Chip. Multimodality and large multimodal models (lmms), 2023. URL \url{https://huyenchip.com/2023/10/10/multimodal.html}.}
\refentry{ch09-ref-011}{A Chowdhery, S Narang, J Devlin, M Bosma, G Mishra, A Roberts, P Barham, HW Chung, and C Sutton. S. gehrmannet al.,“palm: Scalinglanguage modeling with pathways,”. \textit{arXiv preprint arXiv:2204.02311}, 2022.}
\refentry{ch09-ref-012}{Hyung Won Chung, Le Hou, Shayne Longpre, Barret Zoph, Yi Tay, William Fedus, Yunxuan Li, Xuezhi Wang, Mostafa Dehghani, Siddhartha Brahma, et al. Scaling instruction-finetuned language models. \textit{arXiv preprint arXiv:2210.11416}, 2022.}
\refentry{ch09-ref-013}{Chenhang Cui, Yiyang Zhou, Xinyu Yang, Shirley Wu, Linjun Zhang, James Zou, and Huaxiu Yao. Holistic analysis of hallucination in gpt-4v (ision): Bias and interference challenges. \textit{arXiv preprint arXiv:2311.03287}, 2023.}
\refentry{ch09-ref-014}{Alexey Dosovitskiy, Lucas Beyer, Alexander Kolesnikov, Dirk Weissenborn, Xiaohua Zhai, Thomas Unterthiner, Mostafa Dehghani, Matthias Minderer, Georg Heigold, Sylvain Gelly, et al. An image is worth 16x16 words: Transformers for image recognition at scale. \textit{arXiv preprint arXiv:2010.11929}, 2020.}
\refentry{ch09-ref-015}{Danny Driess, Fei Xia, Mehdi SM Sajjadi, Corey Lynch, Aakanksha Chowdhery, Brian Ichter, Ayzaan Wahid, Jonathan Tompson, Quan Vuong, Tianhe Yu, et al. Palm-e: An embodied multimodal language model. \textit{arXiv preprint} \textit{arXiv:2303.03378}, 2023.}
\refentry{ch09-ref-016}{Yuxin Fang, Wen Wang, Binhui Xie, Quan Sun, Ledell Wu, Xinggang Wang, Tiejun Huang, Xinlong Wang, and Yue Cao. Eva: Exploring the limits of masked visual representation learning at scale. In \textit{Proceedings of the IEEE/CVF Conference on} \textit{Computer Vision and Pattern Recognition}, pages 19358–19369, 2023.}
\refentry{ch09-ref-017}{Kirill Gavrilyuk, Ryan Sanford, Mehrsan Javan, and Cees GM Snoek. Actor-transformers for group activity recognition. In \textit{Proceedings of the IEEE/CVF} \textit{Conference on Computer Vision and Pattern Recognition}, pages 839–848, 2020.}
\refentry{ch09-ref-018}{Rohit Girdhar, Alaaeldin El-Nouby, Zhuang Liu, Mannat Singh, Kalyan Vasudev Alwala, Armand Joulin, and Ishan Misra. Imagebind: One embedding space to bind them all. In \textit{Proceedings of the IEEE/CVF Conference on Computer Vision} \textit{and Pattern Recognition}, pages 15180–15190, 2023.}
\refentry{ch09-ref-019}{Tanmay Gupta and Aniruddha Kembhavi. Visual programming: Compositional visual reasoning without training. In \textit{Proceedings of the IEEE/CVF Conference on} \textit{Computer Vision and Pattern Recognition}, pages 14953–14962, 2023.}
\refentry{ch09-ref-020}{Vaishnavi Himakunthala, Andy Ouyang, Daniel Rose, Ryan He, Alex Mei, Yujie Lu, Chinmay Sonar, Michael Saxon, and William Yang Wang. Let’s think frame by frame: Evaluating video chain of thought with video infilling and prediction. \textit{arXiv preprint arXiv:2305.13903}, 2023.}
\refentry{ch09-ref-021}{Geoffrey E Hinton and Russ R Salakhutdinov. A better way to pretrain deep boltzmann machines. \textit{Advances in Neural Information Processing Systems}, 25, 2012.}
\refentry{ch09-ref-022}{Jordan Hoffmann et al. Training compute-optimal large language models, 2022.}
\refentry{ch09-ref-023}{Shaohan Huang, Li Dong, Wenhui Wang, Yaru Hao, Saksham Singhal, Shuming Ma, Tengchao Lv, Lei Cui, Owais Khan Mohammed, Barun Patra, et al. Language is not all you need: Aligning perception with language models. \textit{Advances in Neural} \textit{Information Processing Systems}, 36, 2024.}
\refentry{ch09-ref-024}{Yupan Huang, Zaiqiao Meng, Fangyu Liu, Yixuan Su, Nigel Collier, and Yutong Lu. Sparkles: Unlocking chats across multiple images for multimodal instruction-following models. \textit{arXiv preprint arXiv:2308.16463}, 2023.}
\refentry{ch09-ref-025}{Yuqi Huo, Manli Zhang, Guangzhen Liu, Haoyu Lu, Yizhao Gao, Guoxing Yang, Jingyuan Wen, Heng Zhang, Baogui Xu, Weihao Zheng, et al. Wenlan: Bridging vision and language by large-scale multi-modal pre-training. \textit{arXiv preprint} \textit{arXiv:2103.06561}, 2021.}
\refentry{ch09-ref-026}{Diederik P Kingma and Max Welling. Auto-encoding variational bayes. \textit{arXiv} \textit{preprint arXiv:1312.6114}, 2013.}
\refentry{ch09-ref-027}{COL Stephen A LaRocca, John J Morgan, and Sherri M Bellinger. On the path to 2x learning: Exploring the possibilities of advanced speech recognition. \textit{Calico} \textit{Journal}, pages 295–310, 1999.}
\refentry{ch09-ref-028}{Hugo Laurençon et al. Obelics: An open web-scale filtered dataset of interleaved image-text documents, 2023.}
\refentry{ch09-ref-029}{Gen Li, Nan Duan, Yuejian Fang, Ming Gong, and Daxin Jiang. Unicoder-vl: A universal encoder for vision and language by cross-modal pre-training. In \textit{Proceedings of the AAAI conference on artificial intelligence}, volume 34, pages 11336– 11344, 2020a.}
\refentry{ch09-ref-030}{Lei Li, Zhihui Xie, Mukai Li, Shunian Chen, Peiyi Wang, Liang Chen, Yazheng Yang, Benyou Wang, and Lingpeng Kong. Silkie: Preference distillation for large visual language models. \textit{arXiv preprint arXiv:2312.10665}, 2023.}
\refentry{ch09-ref-031}{Linjie Li, Yen-Chun Chen, Yu Cheng, Zhe Gan, Licheng Yu, and Jingjing Liu. Hero: Hierarchical encoder for video+ language omni-representation pre-training. \textit{arXiv} \textit{preprint arXiv:2005.00200}, 2020b.}
\refentry{ch09-ref-032}{Ruilong Li, Shan Yang, David A Ross, and Angjoo Kanazawa. Ai choreographer: Music conditioned 3d dance generation with aist++. In \textit{Proceedings of the} \textit{IEEE/CVF International Conference on Computer Vision}, pages 13401–13412, 2021.}
\refentry{ch09-ref-033}{Shengzhi Li and Nima Tajbakhsh. Scigraphqa: A large-scale synthetic multi-turn question-answering dataset for scientific graphs. \textit{arXiv preprint} \textit{arXiv:2308.03349}, 2023.}
\refentry{ch09-ref-034}{Xiujun Li, Xi Yin, Chunyuan Li, Pengchuan Zhang, Xiaowei Hu, Lei Zhang, Lijuan Wang, Houdong Hu, Li Dong, Furu Wei, et al. Oscar: Object-semantics aligned pre-training for vision-language tasks. In \textit{Computer Vision–ECCV 2020: 16th} \textit{European Conference, Glasgow, UK, August 23–28, 2020, Proceedings, Part XXX} \textit{16}, pages 121–137. Springer, 2020c.}
\refentry{ch09-ref-035}{Zheng Lian, Licai Sun, Mingyu Xu, Haiyang Sun, Ke Xu, Zhuofan Wen, Shun Chen, Bin Liu, and Jianhua Tao. Explainable multimodal emotion reasoning. \textit{arXiv} \textit{preprint arXiv:2306.15401}, 2023.}
\refentry{ch09-ref-036}{Junyang Lin, An Yang, Yichang Zhang, Jie Liu, Jingren Zhou, and Hongxia Yang. Interbert: Vision-and-language interaction for multi-modal pretraining. \textit{arXiv} \textit{preprint arXiv:2003.13198}, 2020.}
\refentry{ch09-ref-037}{Haohe Liu, Qiao Tian, Yi Yuan, Xubo Liu, Xinhao Mei, Qiuqiang Kong, Yuping Wang, Wenwu Wang, Yuxuan Wang, and Mark D Plumbley. Audioldm 2: Learning holistic audio generation with self-supervised pretraining. \textit{arXiv preprint} \textit{arXiv:2308.05734}, 2023.}
\refentry{ch09-ref-038}{Haotian Liu, Chunyuan Li, Qingyang Wu, and Yong Jae Lee. Visual instruction tuning. \textit{Advances in neural information processing systems}, 36, 2024.}
\refentry{ch09-ref-039}{Jiasen Lu, Dhruv Batra, Devi Parikh, and Stefan Lee. Vilbert: Pretraining task-agnostic visiolinguistic representations for vision-and-language tasks. \textit{Advances} \textit{in neural information processing systems}, 32, 2019.}
\refentry{ch09-ref-040}{Jiquan Ngiam, Aditya Khosla, Mingyu Kim, Juhan Nam, Honglak Lee, and Andrew Y Ng. Multimodal deep learning. In \textit{Proceedings of the 28th international} \textit{conference on machine learning (ICML-11)}, pages 689–696, 2011.}
\refentry{ch09-ref-041}{Munan Ning, Bin Zhu, Yujia Xie, Bin Lin, Jiaxi Cui, Lu Yuan, Dongdong Chen, and Li Yuan. Video-bench: A comprehensive benchmark and toolkit for evaluating video-based large language models. \textit{arXiv preprint arXiv:2311.16103}, 2023.}
\refentry{ch09-ref-042}{Alec Radford, Jong Wook Kim, Chris Hallacy, Aditya Ramesh, Gabriel Goh, Sandhini Agarwal, Girish Sastry, Amanda Askell, Pamela Mishkin, Jack Clark, et al. Learning transferable visual models from natural language supervision. In \textit{International conference on machine learning}, pages 8748–8763. PMLR, 2021.}
\refentry{ch09-ref-043}{Shuhuai Ren, Linli Yao, Shicheng Li, Xu Sun, and Lu Hou. Timechat: A time-sensitive multimodal large language model for long video understanding. \textit{arXiv} \textit{preprint arXiv:2312.02051}, 2023.}
\refentry{ch09-ref-044}{Robin Rombach, Andreas Blattmann, Dominik Lorenz, Patrick Esser, and Björn Ommer. High-resolution image synthesis with latent diffusion models. In \textit{Proceedings of the IEEE/CVF conference on computer vision and pattern recognition}, pages 10684–10695, 2022.}
\refentry{ch09-ref-045}{Olaf Ronneberger, Philipp Fischer, and Thomas Brox. U-net: Convolutional networks for biomedical image segmentation. In \textit{Medical Image Computing and} \textit{Computer-Assisted Intervention–MICCAI 2015: 18th International Conference,} \textit{Munich, Germany, October 5-9, 2015, Proceedings, Part III 18}, pages 234–241. Springer, 2015.}
\refentry{ch09-ref-046}{Shin’ichi Satoh and Takeo Kanade. Name-it: Association of face and name in video. In \textit{Proceedings of IEEE Computer Society Conference on Computer Vision and} \textit{Pattern Recognition}, pages 368–373. IEEE, 1997.}
\refentry{ch09-ref-047}{Bowen Shi, Wei-Ning Hsu, Kushal Lakhotia, and Abdelrahman Mohamed. Learning audio-visual speech representation by masked multimodal cluster prediction. \textit{arXiv preprint arXiv:2201.02184}, 2022.}
\refentry{ch09-ref-048}{Chen Sun, Austin Myers, Carl Vondrick, Kevin Murphy, and Cordelia Schmid. Videobert: A joint model for video and language representation learning. In \textit{Proceedings of the IEEE/CVF international conference on computer vision}, pages 7464–7473, 2019.}
\refentry{ch09-ref-049}{Zhiqing Sun, Sheng Shen, Shengcao Cao, Haotian Liu, Chunyuan Li, Yikang Shen, Chuang Gan, Liang-Yan Gui, Yu-Xiong Wang, Yiming Yang, et al. Aligning large multimodal models with factually augmented rlhf. \textit{arXiv preprint} \textit{arXiv:2309.14525}, 2023.}
\refentry{ch09-ref-050}{Yi Tay, Mostafa Dehghani, Vinh Q Tran, Xavier Garcia, Jason Wei, Xuezhi Wang, Hyung Won Chung, Dara Bahri, Tal Schuster, Steven Zheng, et al. Ul2: Unifying language learning paradigms. In \textit{The Eleventh International Conference on} \textit{Learning Representations}, 2022.}
\refentry{ch09-ref-051}{Shengbang Tong, Zhuang Liu, Yuexiang Zhai, Yi Ma, Yann LeCun, and Saining Xie. Eyes wide shut? exploring the visual shortcomings of multimodal llms. \textit{arXiv} \textit{preprint arXiv:2401.06209}, 2024.}
\refentry{ch09-ref-052}{Hugo Touvron et al. Llama 2: Open foundation and fine-tuned chat models, 2023.}
\refentry{ch09-ref-053}{Gokhan Tur, Andreas Stolcke, Lynn Voss, Stanley Peters, Dilek Hakkani-Tur, John Dowding, Benoit Favre, Raquel Fernández, Matthew Frampton, Mike Frandsen, et al. The calo meeting assistant system. \textit{IEEE Transactions on Audio, Speech,} \textit{and Language Processing}, 18(6):1601–1611, 2010.}
\refentry{ch09-ref-054}{Alessandro Vinciarelli, Maja Pantic, Hervé Bourlard, and Alex Pentland. Social signal processing: state-of-the-art and future perspectives of an emerging domain. In \textit{Proceedings of the 16th ACM international conference on Multimedia}, pages 1061–1070, 2008.}
\refentry{ch09-ref-055}{Xiao Wang, Guangyao Chen, Guangwu Qian, Pengcheng Gao, Xiao-Yong Wei, Yaowei Wang, Yonghong Tian, and Wen Gao. Large-scale multi-modal pre-trained models: A comprehensive survey. \textit{Machine Intelligence Research}, pages 1–36, 2023.}
\refentry{ch09-ref-056}{Zirui Wang, Jiahui Yu, Adams Wei Yu, Zihang Dai, Yulia Tsvetkov, and Yuan Cao. Simvlm: Simple visual language model pretraining with weak supervision. \textit{arXiv} \textit{preprint arXiv:2108.10904}, 2021.}
\refentry{ch09-ref-057}{Chenfei Wu, Shengming Yin, Weizhen Qi, Xiaodong Wang, Zecheng Tang, and Nan Duan. Visual chatgpt: Talking, drawing and editing with visual foundation models. \textit{arXiv preprint arXiv:2303.04671}, 2023a.}
\refentry{ch09-ref-058}{Jiayang Wu, Wensheng Gan, Zefeng Chen, Shicheng Wan, and S Yu Philip. Multimodal large language models: A survey. In \textit{2023 IEEE International Conference} \textit{on Big Data (BigData)}, pages 2247–2256. IEEE, 2023b.}
\refentry{ch09-ref-059}{Shengqiong Wu, Hao Fei, Leigang Qu, Wei Ji, and Tat-Seng Chua. Next-gpt: Any-to-any multimodal llm. \textit{arXiv preprint arXiv:2309.05519}, 2023c.}
\refentry{ch09-ref-060}{Haiyang Xu, Ming Yan, Chenliang Li, Bin Bi, Songfang Huang, Wenming Xiao, and Fei Huang. E2e-vlp: end-to-end vision-language pre-training enhanced by visual learning. \textit{arXiv preprint arXiv:2106.01804}, 2021.}
\refentry{ch09-ref-061}{Peng Xu, Xiatian Zhu, and David A Clifton. Multimodal learning with transformers: A survey. \textit{IEEE Transactions on Pattern Analysis and Machine Intelligence}, 2023.}
\refentry{ch09-ref-062}{Le Xue, Ning Yu, Shu Zhang, Junnan Li, Roberto Martín-Martín, Jiajun Wu, Caiming Xiong, Ran Xu, Juan Carlos Niebles, and Silvio Savarese. Ulip-2: Towards scalable multimodal pre-training for 3d understanding. \textit{arXiv preprint} \textit{arXiv:2305.08275}, 2023.}
\refentry{ch09-ref-063}{Shukang Yin, Chaoyou Fu, Sirui Zhao, Ke Li, Xing Sun, Tong Xu, and Enhong Chen. A survey on multimodal large language models. \textit{arXiv preprint arXiv:2306.13549}, 2023.}
\refentry{ch09-ref-064}{Zhenfei Yin, Jiong Wang, Jianjian Cao, Zhelun Shi, Dingning Liu, Mukai Li, Xiaoshui Huang, Zhiyong Wang, Lu Sheng, Lei Bai, et al. Lamm: Language-assisted multi-modal instruction-tuning dataset, framework, and benchmark. \textit{Advances in} \textit{Neural Information Processing Systems}, 36, 2024.}
\refentry{ch09-ref-065}{Quanzeng You, Hailin Jin, Zhaowen Wang, Chen Fang, and Jiebo Luo. Image captioning with semantic attention. In \textit{Proceedings of the IEEE conference on computer vision and pattern recognition}, pages 4651–4659, 2016.}
\refentry{ch09-ref-066}{Peter Young, Alice Lai, Micah Hodosh, and Julia Hockenmaier. From image descriptions to visual denotations: New similarity metrics for semantic inference over event descriptions. \textit{Transactions of the Association for Computational Linguistics}, 2:67–78, 2014. doi: 10.1162/tacl\_a\_00166. URL \url{https://aclanthology.org/Q14-1006}.}
\refentry{ch09-ref-067}{Tianyu Yu, Yuan Yao, Haoye Zhang, Taiwen He, Yifeng Han, Ganqu Cui, Jinyi Hu, Zhiyuan Liu, Hai-Tao Zheng, Maosong Sun, et al. Rlhf-v: Towards trustworthy mllms via behavior alignment from fine-grained correctional human feedback. \textit{arXiv preprint arXiv:2312.00849}, 2023.}
\refentry{ch09-ref-068}{Aohan Zeng, Xiao Liu, Zhengxiao Du, Zihan Wang, Hanyu Lai, Ming Ding, Zhuoyi Yang, Yifan Xu, Wendi Zheng, Xiao Xia, et al. Glm-130b: An open bilingual pre-trained model. \textit{arXiv preprint arXiv:2210.02414}, 2022.}
\refentry{ch09-ref-069}{Xunlin Zhan, Yangxin Wu, Xiao Dong, Yunchao Wei, Minlong Lu, Yichi Zhang, Hang Xu, and Xiaodan Liang. Product1m: Towards weakly supervised instance-level product retrieval via cross-modal pretraining. In \textit{Proceedings of the} \textit{IEEE/CVF International Conference on Computer Vision}, pages 11782–11791, 2021.}
\refentry{ch09-ref-070}{Duzhen Zhang, Yahan Yu, Chenxing Li, Jiahua Dong, Dan Su, Chenhui Chu, and Dong Yu. Mm-llms: Recent advances in multimodal large language models. \textit{arXiv} \textit{preprint arXiv:2401.13601}, 2024a.}
\refentry{ch09-ref-071}{Ge Zhang, Xinrun Du, Bei Chen, Yiming Liang, Tongxu Luo, Tianyu Zheng, Kang Zhu, Yuyang Cheng, Chunpu Xu, Shuyue Guo, et al. Cmmmu: A chinese massive multi-discipline multimodal understanding benchmark. \textit{arXiv preprint} \textit{arXiv:2401.11944}, 2024b.}
\refentry{ch09-ref-072}{Hang Zhang, Xin Li, and Lidong Bing. Video-llama: An instruction-tuned audio-visual language model for video understanding. \textit{arXiv preprint arXiv:2306.02858}, 2023.}
\refentry{ch09-ref-073}{Jiawei Zhang, Tianyu Pang, Chao Du, Yi Ren, Bo Li, and Min Lin. Benchmarking large multimodal models against common corruptions. \textit{arXiv preprint} \textit{arXiv:2401.11943}, 2024c.}
\refentry{ch09-ref-074}{Susan Zhang, Stephen Roller, Naman Goyal, Mikel Artetxe, Moya Chen, Shuohui Chen, Christopher Dewan, Mona Diab, Xian Li, Xi Victoria Lin, et al. Opt: Open pre-trained transformer language models. \textit{arXiv preprint arXiv:2205.01068}, 2022.}
\refentry{ch09-ref-075}{Min Zhao, Fan Bao, Chongxuan Li, and Jun Zhu. Egsde: Unpaired image-to-image translation via energy-guided stochastic differential equations. \textit{Advances in Neural} \textit{Information Processing Systems}, 35:3609–3623, 2022.}
\refentry{ch09-ref-076}{Zijia Zhao, Longteng Guo, Tongtian Yue, Sihan Chen, Shuai Shao, Xinxin Zhu, Zehuan Yuan, and Jing Liu. Chatbridge: Bridging modalities with large language model as a language catalyst. \textit{arXiv preprint arXiv:2305.16103}, 2023.}
\refentry{ch09-ref-077}{Deyao Zhu, Jun Chen, Xiaoqian Shen, Xiang Li, and Mohamed Elhoseiny. Minigpt-4: Enhancing vision-language understanding with advanced large language models. \textit{arXiv preprint arXiv:2304.10592}, 2023.}
\end{chapterreferences}

